%%%%%%%%%%%%%%%%%%%%%%%%%%%%%%%%%%%%%%%%%%%%%%%%%%%%%%%%%%%%
% 11.9 COMPUTER ALGEBRA
% The Composition of Advanced Mathematics
%%%%%%%%%%%%%%%%%%%%%%%%%%%%%%%%%%%%%%%%%%%%%%%%%%%%%%%%%%%%

\section{Computer Algebra}
\label{sec:computer-algebra}

\prerequisite{
Abstract algebra, polynomial arithmetic, rings and fields, ideals,
elementary number theory, symbolic computation, algorithms, data
structures, and basic computational complexity.
}

\learningobjectives{
By the end of this section, the reader should be able to:
\begin{itemize}
    \item explain the goals of computer algebra;
    \item distinguish symbolic computation from computer algebra;
    \item understand algebraic data representations;
    \item represent integers, rationals, polynomials, and rational
          functions exactly;
    \item distinguish dense and sparse polynomial representations;
    \item understand coefficient domains;
    \item recognize polynomial rings;
    \item understand multivariate monomials;
    \item work with monomial orderings;
    \item recognize lexicographic and graded orderings;
    \item perform polynomial addition and multiplication algorithmically;
    \item understand polynomial division;
    \item recognize pseudo-division over nonfields;
    \item compute polynomial greatest common divisors;
    \item understand subresultant ideas conceptually;
    \item recognize square-free factorization;
    \item understand polynomial factorization over finite fields;
    \item recognize rational polynomial factorization;
    \item understand modular factorization strategies;
    \item recognize Hensel lifting conceptually;
    \item understand the Chinese Remainder Theorem in computer algebra;
    \item understand rational reconstruction;
    \item recognize modular algorithms for exact computation;
    \item understand algebraic numbers;
    \item represent algebraic extensions;
    \item recognize minimal polynomials;
    \item perform arithmetic modulo polynomial relations;
    \item understand polynomial ideals computationally;
    \item define leading terms and leading monomials;
    \item understand multivariate polynomial reduction;
    \item define Gröbner bases;
    \item recognize Buchberger's algorithm;
    \item understand S-polynomials;
    \item understand ideal membership testing;
    \item recognize elimination using Gröbner bases;
    \item understand solving zero-dimensional polynomial systems
          conceptually;
    \item recognize resultants and Sylvester matrices;
    \item understand elimination theory;
    \item recognize polynomial system solving over exact fields;
    \item distinguish exact and numerical algebraic geometry;
    \item understand symbolic linear algebra;
    \item recognize fraction-free elimination;
    \item understand determinant algorithms;
    \item recognize exact matrix normal forms;
    \item understand Hermite and Smith normal forms conceptually;
    \item recognize sparse exact linear algebra;
    \item understand expression swell;
    \item recognize modular methods for matrix computations;
    \item understand fast integer arithmetic conceptually;
    \item recognize fast polynomial multiplication;
    \item understand FFT-based polynomial multiplication conceptually;
    \item recognize computational complexity in algebraic algorithms;
    \item understand bit complexity;
    \item recognize coefficient growth;
    \item understand probabilistic and randomized algebra algorithms;
    \item distinguish Las Vegas and Monte Carlo algebra algorithms;
    \item understand computer algebra system architecture conceptually;
    \item recognize kernels, simplifiers, domains, and evaluators;
    \item understand canonicalization;
    \item recognize pattern matching and rewriting within computer algebra;
    \item understand symbolic-numeric hybrid algorithms;
    \item connect computer algebra with algebra, number theory, algebraic
          geometry, cryptography, coding theory, optimization, and
          scientific computing.
\end{itemize}
}

%%%%%%%%%%%%%%%%%%%%%%%%%%%%%%%%%%%%%%%%%%%%%%%%%%%%%%%%%%%%
\subsection{What Is Computer Algebra?}
\label{subsec:what-is-computer-algebra}
%%%%%%%%%%%%%%%%%%%%%%%%%%%%%%%%%%%%%%%%%%%%%%%%%%%%%%%%%%%%

Computer algebra studies algorithms and data structures for exact
mathematical computation.

Its central objects include:

\[
\boxed{
\text{integers},
}
\]

\[
\boxed{
\text{rational numbers},
}
\]

\[
\boxed{
\text{polynomials},
}
\]

\[
\boxed{
\text{algebraic numbers},
}
\]

\[
\boxed{
\text{matrices},
}
\]

and

\[
\boxed{
\text{symbolic expressions}.
}
\]

The emphasis is on exact algebraic manipulation rather than approximate
floating-point computation.

%%%%%%%%%%%%%%%%%%%%%%%%%%%%%%%%%%%%%%%%%%%%%%%%%%%%%%%%%%%%
\subsection{Computer Algebra Versus Symbolic Computation}
\label{subsec:computer-algebra-vs-symbolic}
%%%%%%%%%%%%%%%%%%%%%%%%%%%%%%%%%%%%%%%%%%%%%%%%%%%%%%%%%%%%

Symbolic computation is the broad practice of manipulating mathematical
expressions.

Computer algebra focuses more specifically on the algorithms that make
exact algebraic manipulation possible.

Thus:

\[
\boxed{
\text{symbolic computation}
=
\text{user-level exact manipulation},
}
\]

while

\[
\boxed{
\text{computer algebra}
=
\text{algorithmic machinery behind exact manipulation}.
}
\]

The two subjects overlap heavily.

%%%%%%%%%%%%%%%%%%%%%%%%%%%%%%%%%%%%%%%%%%%%%%%%%%%%%%%%%%%%
\subsection{Exact Algebraic Domains}
\label{subsec:exact-algebraic-domains}
%%%%%%%%%%%%%%%%%%%%%%%%%%%%%%%%%%%%%%%%%%%%%%%%%%%%%%%%%%%%

A computer algebra system must know the algebraic domain in which an object
lives.

Examples include:

\[
\boxed{
\mathbb{Z},
}
\]

\[
\boxed{
\mathbb{Q},
}
\]

\[
\boxed{
\mathbb{F}_p,
}
\]

\[
\boxed{
\mathbb{Q}[x],
}
\]

and

\[
\boxed{
\mathbb{Q}(x).
}
\]

The domain determines which operations are permitted and which algorithms
apply.

%%%%%%%%%%%%%%%%%%%%%%%%%%%%%%%%%%%%%%%%%%%%%%%%%%%%%%%%%%%%
\subsection{Integers}
\label{subsec:computer-algebra-integers}
%%%%%%%%%%%%%%%%%%%%%%%%%%%%%%%%%%%%%%%%%%%%%%%%%%%%%%%%%%%%

Arbitrary-size integers are fundamental exact objects.

Unlike fixed-width machine integers, their size is limited primarily by
available memory.

For

\[
a,b\in\mathbb{Z},
\]

a computer algebra system must support exact:

\[
\boxed{
a+b,
}
\]

\[
\boxed{
ab,
}
\]

\[
\boxed{
\gcd(a,b),
}
\]

and

\[
\boxed{
a\bmod b.
}
\]

Large-integer arithmetic is the foundation of many higher-level symbolic
algorithms.

%%%%%%%%%%%%%%%%%%%%%%%%%%%%%%%%%%%%%%%%%%%%%%%%%%%%%%%%%%%%
\subsection{Rational Numbers}
\label{subsec:computer-algebra-rationals}
%%%%%%%%%%%%%%%%%%%%%%%%%%%%%%%%%%%%%%%%%%%%%%%%%%%%%%%%%%%%

A rational number is usually stored as

\[
\boxed{
\frac{p}{q},
}
\]

where

\[
p,q\in\mathbb{Z},
\qquad
q>0,
\]

and

\[
\boxed{
\gcd(p,q)=1.
}
\]

This normalization gives a canonical representation.

For example,

\[
\frac{-8}{-12}
\]

becomes

\[
\boxed{
\frac23.
}
\]

%%%%%%%%%%%%%%%%%%%%%%%%%%%%%%%%%%%%%%%%%%%%%%%%%%%%%%%%%%%%
\subsection{Normalization of Rational Arithmetic}
\label{subsec:rational-normalization}
%%%%%%%%%%%%%%%%%%%%%%%%%%%%%%%%%%%%%%%%%%%%%%%%%%%%%%%%%%%%

Suppose

\[
\frac{a}{b}
+
\frac{c}{d}
=
\frac{
ad+bc
}{
bd
}.
\]

The numerator and denominator may share a common factor.

A normalized result divides by

\[
\boxed{
g
=
\gcd(ad+bc,bd).
}
\]

Repeated normalization prevents uncontrolled growth of equivalent
fraction representations.

%%%%%%%%%%%%%%%%%%%%%%%%%%%%%%%%%%%%%%%%%%%%%%%%%%%%%%%%%%%%
\subsection{Polynomial Rings}
\label{subsec:computer-algebra-polynomial-rings}
%%%%%%%%%%%%%%%%%%%%%%%%%%%%%%%%%%%%%%%%%%%%%%%%%%%%%%%%%%%%

A polynomial ring in variable \(x\) over coefficient ring \(R\) is denoted

\[
\boxed{
R[x].
}
\]

For example,

\[
\mathbb{Q}[x]
\]

contains polynomials

\[
\boxed{
a_0+a_1x+\cdots+a_nx^n,
\qquad
a_i\in\mathbb{Q}.
}
\]

For several variables,

\[
\boxed{
R[x_1,\ldots,x_n].
}
\]

The coefficient domain plays a central role in algorithm selection.

%%%%%%%%%%%%%%%%%%%%%%%%%%%%%%%%%%%%%%%%%%%%%%%%%%%%%%%%%%%%
\subsection{Dense Polynomial Representation}
\label{subsec:dense-polynomial-representation}
%%%%%%%%%%%%%%%%%%%%%%%%%%%%%%%%%%%%%%%%%%%%%%%%%%%%%%%%%%%%

A dense univariate polynomial

\[
p(x)
=
a_0+a_1x+\cdots+a_nx^n
\]

may be stored as

\[
\boxed{
[a_0,a_1,\ldots,a_n].
}
\]

This is efficient when many coefficients are nonzero.

For example,

\[
2-3x+5x^2
\]

may be stored as

\[
\boxed{
[2,-3,5].
}
\]

%%%%%%%%%%%%%%%%%%%%%%%%%%%%%%%%%%%%%%%%%%%%%%%%%%%%%%%%%%%%
\subsection{Sparse Polynomial Representation}
\label{subsec:sparse-polynomial-representation}
%%%%%%%%%%%%%%%%%%%%%%%%%%%%%%%%%%%%%%%%%%%%%%%%%%%%%%%%%%%%

A sparse polynomial stores only nonzero terms.

For example,

\[
x^{1000}+3x^7-2
\]

may be stored conceptually as

\[
\boxed{
\{
(1000,1),
(7,3),
(0,-2)
\}.
}
\]

This avoids storing hundreds of zero coefficients.

%%%%%%%%%%%%%%%%%%%%%%%%%%%%%%%%%%%%%%%%%%%%%%%%%%%%%%%%%%%%
\subsection{Dense Versus Sparse Representation}
\label{subsec:dense-vs-sparse-polynomial}
%%%%%%%%%%%%%%%%%%%%%%%%%%%%%%%%%%%%%%%%%%%%%%%%%%%%%%%%%%%%

Dense representation is efficient for moderate-degree polynomials with
many nonzero coefficients.

Sparse representation is efficient when

\[
\boxed{
\#\text{nonzero terms}
\ll
\deg p.
}
\]

The best representation depends on the algebraic structure of the
problem.

%%%%%%%%%%%%%%%%%%%%%%%%%%%%%%%%%%%%%%%%%%%%%%%%%%%%%%%%%%%%
\subsection{Multivariate Monomials}
\label{subsec:multivariate-monomials}
%%%%%%%%%%%%%%%%%%%%%%%%%%%%%%%%%%%%%%%%%%%%%%%%%%%%%%%%%%%%

A multivariate monomial has form

\[
\boxed{
x_1^{\alpha_1}
x_2^{\alpha_2}
\cdots
x_n^{\alpha_n},
}
\]

where

\[
\alpha_i\in\mathbb{Z}_{\geq0}.
\]

The exponent vector is

\[
\boxed{
\alpha
=
(\alpha_1,\ldots,\alpha_n).
}
\]

Thus a polynomial can be represented as a mapping from exponent vectors to
coefficients.

%%%%%%%%%%%%%%%%%%%%%%%%%%%%%%%%%%%%%%%%%%%%%%%%%%%%%%%%%%%%
\subsection{Monomial Ordering}
\label{subsec:monomial-ordering}
%%%%%%%%%%%%%%%%%%%%%%%%%%%%%%%%%%%%%%%%%%%%%%%%%%%%%%%%%%%%

Multivariate algorithms require an ordering on monomials.

A monomial order must be:

\[
\boxed{
\text{total},
}
\]

\[
\boxed{
\text{multiplicatively compatible},
}
\]

and

\[
\boxed{
\text{well ordered}.
}
\]

Different orderings produce different Gröbner bases and different
computational behavior.

%%%%%%%%%%%%%%%%%%%%%%%%%%%%%%%%%%%%%%%%%%%%%%%%%%%%%%%%%%%%
\subsection{Lexicographic Order}
\label{subsec:lexicographic-order}
%%%%%%%%%%%%%%%%%%%%%%%%%%%%%%%%%%%%%%%%%%%%%%%%%%%%%%%%%%%%

In lexicographic order, compare exponent vectors from the first variable
onward.

For

\[
x>y,
\]

we have

\[
\boxed{
x^2
>
xy^{100}
}
\]

because the exponent of \(x\) is compared first.

Lexicographic order is particularly useful for elimination.

%%%%%%%%%%%%%%%%%%%%%%%%%%%%%%%%%%%%%%%%%%%%%%%%%%%%%%%%%%%%
\subsection{Graded Lexicographic Order}
\label{subsec:graded-lexicographic-order}
%%%%%%%%%%%%%%%%%%%%%%%%%%%%%%%%%%%%%%%%%%%%%%%%%%%%%%%%%%%%

Graded lexicographic order first compares total degree.

For monomials

\[
x^\alpha
\]

and

\[
x^\beta,
\]

compare

\[
|\alpha|
=
\sum_i\alpha_i
\]

and

\[
|\beta|.
\]

If total degrees tie, lexicographic order breaks the tie.

%%%%%%%%%%%%%%%%%%%%%%%%%%%%%%%%%%%%%%%%%%%%%%%%%%%%%%%%%%%%
\subsection{Graded Reverse Lexicographic Order}
\label{subsec:grevlex}
%%%%%%%%%%%%%%%%%%%%%%%%%%%%%%%%%%%%%%%%%%%%%%%%%%%%%%%%%%%%

Graded reverse lexicographic order also compares total degree first.

Ties are resolved by examining exponent differences from the last
variable backward.

This ordering often performs efficiently in Gröbner-basis computations.

%%%%%%%%%%%%%%%%%%%%%%%%%%%%%%%%%%%%%%%%%%%%%%%%%%%%%%%%%%%%
\subsection{Leading Monomial and Leading Term}
\label{subsec:leading-monomial-term}
%%%%%%%%%%%%%%%%%%%%%%%%%%%%%%%%%%%%%%%%%%%%%%%%%%%%%%%%%%%%

Given a chosen monomial order, a nonzero polynomial has a largest
monomial.

The leading monomial is

\[
\boxed{
\operatorname{LM}(f).
}
\]

If its coefficient is included, the leading term is

\[
\boxed{
\operatorname{LT}(f).
}
\]

For example, under lexicographic order with \(x>y\),

\[
f=x^2+xy^5+y^{20}
\]

has

\[
\boxed{
\operatorname{LM}(f)=x^2.
}
\]

%%%%%%%%%%%%%%%%%%%%%%%%%%%%%%%%%%%%%%%%%%%%%%%%%%%%%%%%%%%%
\subsection{Polynomial Addition Algorithm}
\label{subsec:polynomial-addition-algorithm}
%%%%%%%%%%%%%%%%%%%%%%%%%%%%%%%%%%%%%%%%%%%%%%%%%%%%%%%%%%%%

Polynomial addition combines coefficients of matching monomials.

For sparse representations, terms can be merged according to monomial
order.

If

\[
f
=
3x^2+2x+1
\]

and

\[
g
=
x^2-5x+4,
\]

then

\[
\boxed{
f+g
=
4x^2-3x+5.
}
\]

%%%%%%%%%%%%%%%%%%%%%%%%%%%%%%%%%%%%%%%%%%%%%%%%%%%%%%%%%%%%
\subsection{Polynomial Multiplication Algorithm}
\label{subsec:polynomial-multiplication-algorithm}
%%%%%%%%%%%%%%%%%%%%%%%%%%%%%%%%%%%%%%%%%%%%%%%%%%%%%%%%%%%%

For

\[
f(x)
=
\sum_i a_ix^i
\]

and

\[
g(x)
=
\sum_j b_jx^j,
\]

the coefficient of \(x^k\) in the product is

\[
\boxed{
c_k
=
\sum_{i+j=k}
a_ib_j.
}
\]

The naive dense multiplication algorithm requires approximately

\[
\boxed{
O(n^2)
}
\]

coefficient operations for degree-\(n\) inputs of comparable size.

%%%%%%%%%%%%%%%%%%%%%%%%%%%%%%%%%%%%%%%%%%%%%%%%%%%%%%%%%%%%
\subsection{Fast Polynomial Multiplication}
\label{subsec:fast-polynomial-multiplication}
%%%%%%%%%%%%%%%%%%%%%%%%%%%%%%%%%%%%%%%%%%%%%%%%%%%%%%%%%%%%

Polynomial multiplication can be accelerated using divide-and-conquer or
transform methods.

Evaluation-based methods use the principle

\[
\boxed{
\text{evaluate}
\rightarrow
\text{multiply values}
\rightarrow
\text{interpolate}.
}
\]

FFT-based multiplication can reduce complexity toward

\[
\boxed{
O(n\log n)
}
\]

up to logarithmic and coefficient-arithmetic factors in suitable settings.

%%%%%%%%%%%%%%%%%%%%%%%%%%%%%%%%%%%%%%%%%%%%%%%%%%%%%%%%%%%%
\subsection{Polynomial Division Over a Field}
\label{subsec:polynomial-division-field}
%%%%%%%%%%%%%%%%%%%%%%%%%%%%%%%%%%%%%%%%%%%%%%%%%%%%%%%%%%%%

For

\[
f,g\in F[x],
\qquad
g\neq0,
\]

where \(F\) is a field, there exist unique \(q,r\) satisfying

\[
\boxed{
f=qg+r,
}
\]

with

\[
\boxed{
\deg r<\deg g.
}
\]

Division works because the leading coefficient of \(g\) is invertible in
the field.

%%%%%%%%%%%%%%%%%%%%%%%%%%%%%%%%%%%%%%%%%%%%%%%%%%%%%%%%%%%%
\subsection{Worked Example: Polynomial Division}
\label{subsec:worked-polynomial-division}
%%%%%%%%%%%%%%%%%%%%%%%%%%%%%%%%%%%%%%%%%%%%%%%%%%%%%%%%%%%%

\begin{workedexamplebox}{Exact Polynomial Division}

Divide

\[
x^3-1
\]

by

\[
x-1.
\]

Since

\[
x^3-1
=
(x-1)(x^2+x+1),
\]

we obtain

\begin{answerbox}
\[
\boxed{
q(x)
=
x^2+x+1,
\qquad
r(x)=0.
}
\]
\end{answerbox}

\end{workedexamplebox}

%%%%%%%%%%%%%%%%%%%%%%%%%%%%%%%%%%%%%%%%%%%%%%%%%%%%%%%%%%%%
\subsection{Pseudo-Division}
\label{subsec:pseudo-division}
%%%%%%%%%%%%%%%%%%%%%%%%%%%%%%%%%%%%%%%%%%%%%%%%%%%%%%%%%%%%

Over a coefficient ring where leading coefficients may not be invertible,
ordinary polynomial division may introduce fractions.

Pseudo-division avoids these fractions by multiplying the dividend by a
power of the divisor's leading coefficient.

Conceptually,

\[
\boxed{
c^k f
=
qg+r,
}
\]

where \(c\) is the leading coefficient of \(g\).

This is useful in exact polynomial gcd algorithms over \(\mathbb{Z}[x]\).

%%%%%%%%%%%%%%%%%%%%%%%%%%%%%%%%%%%%%%%%%%%%%%%%%%%%%%%%%%%%
\subsection{Polynomial GCD}
\label{subsec:computer-algebra-polynomial-gcd}
%%%%%%%%%%%%%%%%%%%%%%%%%%%%%%%%%%%%%%%%%%%%%%%%%%%%%%%%%%%%

For

\[
f,g\in F[x],
\]

their gcd can be computed using the Euclidean algorithm.

Repeatedly compute

\[
\boxed{
r_{k+1}
=
r_{k-1}
\bmod
r_k.
}
\]

The last nonzero remainder is the gcd up to multiplication by a nonzero
constant.

The result is usually normalized to be monic.

%%%%%%%%%%%%%%%%%%%%%%%%%%%%%%%%%%%%%%%%%%%%%%%%%%%%%%%%%%%%
\subsection{Worked Example: Euclidean Algorithm}
\label{subsec:worked-polynomial-euclidean}
%%%%%%%%%%%%%%%%%%%%%%%%%%%%%%%%%%%%%%%%%%%%%%%%%%%%%%%%%%%%

\begin{workedexamplebox}{Polynomial GCD}

Let

\[
f=x^3-x
\]

and

\[
g=x^2-1.
\]

Since

\[
x^3-x
=
x(x^2-1),
\]

the remainder after dividing \(f\) by \(g\) is zero.

Therefore,

\begin{answerbox}
\[
\boxed{
\gcd(f,g)
=
x^2-1.
}
\]
\end{answerbox}

\end{workedexamplebox}

%%%%%%%%%%%%%%%%%%%%%%%%%%%%%%%%%%%%%%%%%%%%%%%%%%%%%%%%%%%%
\subsection{Primitive Part and Content}
\label{subsec:primitive-part-content}
%%%%%%%%%%%%%%%%%%%%%%%%%%%%%%%%%%%%%%%%%%%%%%%%%%%%%%%%%%%%

For

\[
f(x)
=
a_0+a_1x+\cdots+a_nx^n
\in
\mathbb{Z}[x],
\]

define the content

\[
\boxed{
\operatorname{cont}(f)
=
\gcd(a_0,\ldots,a_n).
}
\]

Then write

\[
\boxed{
f
=
\operatorname{cont}(f)
\operatorname{pp}(f),
}
\]

where \(\operatorname{pp}(f)\) is the primitive part.

This separation is useful in gcd and factorization algorithms.

%%%%%%%%%%%%%%%%%%%%%%%%%%%%%%%%%%%%%%%%%%%%%%%%%%%%%%%%%%%%
\subsection{Gauss's Lemma Computationally}
\label{subsec:gauss-lemma-computational}
%%%%%%%%%%%%%%%%%%%%%%%%%%%%%%%%%%%%%%%%%%%%%%%%%%%%%%%%%%%%

Primitive polynomial arithmetic over

\[
\mathbb{Z}[x]
\]

is closely related to arithmetic over

\[
\mathbb{Q}[x].
\]

Gauss's lemma allows factorization and gcd questions over the rationals to
be reduced to primitive integer polynomials.

This avoids unnecessary rational coefficient growth.

%%%%%%%%%%%%%%%%%%%%%%%%%%%%%%%%%%%%%%%%%%%%%%%%%%%%%%%%%%%%
\subsection{Subresultant Polynomial Remainder Sequences}
\label{subsec:subresultant-prs}
%%%%%%%%%%%%%%%%%%%%%%%%%%%%%%%%%%%%%%%%%%%%%%%%%%%%%%%%%%%%

Naive pseudo-remainder sequences can produce very large coefficients.

Subresultant polynomial remainder sequences scale intermediate remainders
to control coefficient growth.

They are important in exact:

\[
\boxed{
\text{gcd computation},
}
\]

\[
\boxed{
\text{resultants},
}
\]

and

\[
\boxed{
\text{real-root algorithms}.
}
\]

%%%%%%%%%%%%%%%%%%%%%%%%%%%%%%%%%%%%%%%%%%%%%%%%%%%%%%%%%%%%
\subsection{Square-Free Factorization}
\label{subsec:computer-algebra-square-free}
%%%%%%%%%%%%%%%%%%%%%%%%%%%%%%%%%%%%%%%%%%%%%%%%%%%%%%%%%%%%

Suppose

\[
f\in F[x]
\]

over characteristic zero.

Repeated factors can be detected using

\[
\boxed{
g
=
\gcd(f,f').
}
\]

If

\[
g=1,
\]

then \(f\) is square-free.

If \(g\neq1\), repeated factors are present.

Square-free decomposition separates factors according to multiplicity.

%%%%%%%%%%%%%%%%%%%%%%%%%%%%%%%%%%%%%%%%%%%%%%%%%%%%%%%%%%%%
\subsection{Worked Example: Repeated Factor Detection}
\label{subsec:worked-square-free}
%%%%%%%%%%%%%%%%%%%%%%%%%%%%%%%%%%%%%%%%%%%%%%%%%%%%%%%%%%%%

\begin{workedexamplebox}{Detecting Multiplicity}

Let

\[
f(x)
=
(x-1)^3(x+2).
\]

Then

\[
f'
\]

also contains a factor

\[
(x-1)^2.
\]

Therefore,

\[
\gcd(f,f')
\neq1.
\]

\begin{answerbox}
\[
\boxed{
f
\text{ is not square-free.}
}
\]
\end{answerbox}

The common factor reveals repeated roots.

\end{workedexamplebox}

%%%%%%%%%%%%%%%%%%%%%%%%%%%%%%%%%%%%%%%%%%%%%%%%%%%%%%%%%%%%
\subsection{Factorization Over Finite Fields}
\label{subsec:finite-field-factorization}
%%%%%%%%%%%%%%%%%%%%%%%%%%%%%%%%%%%%%%%%%%%%%%%%%%%%%%%%%%%%

Consider

\[
f(x)
\in
\mathbb{F}_p[x].
\]

Finite-field factorization algorithms typically break the task into:

\[
\boxed{
\text{square-free factorization},
}
\]

\[
\boxed{
\text{distinct-degree factorization},
}
\]

and

\[
\boxed{
\text{equal-degree factorization}.
}
\]

Finite fields provide a computationally convenient environment because
there are only finitely many coefficients.

%%%%%%%%%%%%%%%%%%%%%%%%%%%%%%%%%%%%%%%%%%%%%%%%%%%%%%%%%%%%
\subsection{Why Finite Fields Matter}
\label{subsec:why-finite-fields-computer-algebra}
%%%%%%%%%%%%%%%%%%%%%%%%%%%%%%%%%%%%%%%%%%%%%%%%%%%%%%%%%%%%

Finite-field computations appear in:

\[
\boxed{
\text{coding theory},
}
\]

\[
\boxed{
\text{cryptography},
}
\]

\[
\boxed{
\text{polynomial factorization},
}
\]

and

\[
\boxed{
\text{modular algorithms}.
}
\]

They also provide efficient intermediate domains for exact calculations
over the integers and rationals.

%%%%%%%%%%%%%%%%%%%%%%%%%%%%%%%%%%%%%%%%%%%%%%%%%%%%%%%%%%%%
\subsection{Modular Polynomial Factorization}
\label{subsec:modular-polynomial-factorization}
%%%%%%%%%%%%%%%%%%%%%%%%%%%%%%%%%%%%%%%%%%%%%%%%%%%%%%%%%%%%

To factor a polynomial over

\[
\mathbb{Z},
\]

one may reduce it modulo a suitable prime \(p\):

\[
\boxed{
f(x)
\mapsto
f(x)\bmod p.
}
\]

Factor the smaller polynomial in

\[
\mathbb{F}_p[x].
\]

Then lift and reconstruct factors over the integers.

This avoids large intermediate integer computations.

%%%%%%%%%%%%%%%%%%%%%%%%%%%%%%%%%%%%%%%%%%%%%%%%%%%%%%%%%%%%
\subsection{Hensel Lifting Preview}
\label{subsec:hensel-lifting}
%%%%%%%%%%%%%%%%%%%%%%%%%%%%%%%%%%%%%%%%%%%%%%%%%%%%%%%%%%%%

Suppose a factorization is known modulo \(p\):

\[
\boxed{
f
\equiv
gh
\pmod p.
}
\]

Under suitable coprimality conditions, Hensel lifting improves this to a
factorization modulo

\[
\boxed{
p^2,
p^4,
\ldots
}
\]

or successive powers of \(p\).

Eventually enough modular information is available to reconstruct integer
factors.

%%%%%%%%%%%%%%%%%%%%%%%%%%%%%%%%%%%%%%%%%%%%%%%%%%%%%%%%%%%%
\subsection{Chinese Remainder Reconstruction}
\label{subsec:chinese-remainder-reconstruction}
%%%%%%%%%%%%%%%%%%%%%%%%%%%%%%%%%%%%%%%%%%%%%%%%%%%%%%%%%%%%

Suppose an integer \(a\) is known modulo several pairwise coprime moduli:

\[
a\equiv a_i\pmod{m_i}.
\]

The Chinese Remainder Theorem reconstructs \(a\) modulo

\[
\boxed{
M
=
\prod_i m_i.
}
\]

If the true integer is known to lie within a sufficiently small range, the
exact value can be recovered.

%%%%%%%%%%%%%%%%%%%%%%%%%%%%%%%%%%%%%%%%%%%%%%%%%%%%%%%%%%%%
\subsection{Worked Example: Chinese Remainder Reconstruction}
\label{subsec:worked-chinese-remainder}
%%%%%%%%%%%%%%%%%%%%%%%%%%%%%%%%%%%%%%%%%%%%%%%%%%%%%%%%%%%%

\begin{workedexamplebox}{Reconstructing an Integer}

Suppose

\[
x\equiv2\pmod3,
\]

and

\[
x\equiv3\pmod5.
\]

Numbers congruent to \(2\) modulo \(3\) are

\[
2,5,8,11,14,\ldots.
\]

Among these,

\[
8\equiv3\pmod5.
\]

Therefore,

\begin{answerbox}
\[
\boxed{
x\equiv8\pmod{15}.
}
\]
\end{answerbox}

\end{workedexamplebox}

%%%%%%%%%%%%%%%%%%%%%%%%%%%%%%%%%%%%%%%%%%%%%%%%%%%%%%%%%%%%
\subsection{Rational Reconstruction}
\label{subsec:rational-reconstruction}
%%%%%%%%%%%%%%%%%%%%%%%%%%%%%%%%%%%%%%%%%%%%%%%%%%%%%%%%%%%%

Suppose modular computation gives

\[
a
\pmod M.
\]

One may seek integers \(p,q\) satisfying

\[
\boxed{
a
\equiv
pq^{-1}
\pmod M
}
\]

with \(p\) and \(q\) sufficiently small.

Under appropriate bounds, the rational number

\[
\boxed{
\frac pq
}
\]

can be recovered uniquely.

This is useful when exact rational coefficients are reconstructed from
modular images.

%%%%%%%%%%%%%%%%%%%%%%%%%%%%%%%%%%%%%%%%%%%%%%%%%%%%%%%%%%%%
\subsection{Modular Algorithms}
\label{subsec:modular-algorithms}
%%%%%%%%%%%%%%%%%%%%%%%%%%%%%%%%%%%%%%%%%%%%%%%%%%%%%%%%%%%%

A general modular strategy is:

\[
\boxed{
\text{large exact problem}
}
\]

\[
\downarrow
\]

\[
\boxed{
\text{many small modular problems}
}
\]

\[
\downarrow
\]

\[
\boxed{
\text{reconstruction of exact result}.
}
\]

This technique reduces coefficient growth and enables parallel
computation.

%%%%%%%%%%%%%%%%%%%%%%%%%%%%%%%%%%%%%%%%%%%%%%%%%%%%%%%%%%%%
\subsection{Algebraic Numbers}
\label{subsec:computer-algebra-algebraic-numbers}
%%%%%%%%%%%%%%%%%%%%%%%%%%%%%%%%%%%%%%%%%%%%%%%%%%%%%%%%%%%%

A complex number \(\alpha\) is algebraic over \(\mathbb{Q}\) if it is a
root of a nonzero polynomial

\[
\boxed{
p(x)\in\mathbb{Q}[x].
}
\]

For example,

\[
\sqrt{2}
\]

is algebraic because

\[
\boxed{
x^2-2=0.
}
\]

%%%%%%%%%%%%%%%%%%%%%%%%%%%%%%%%%%%%%%%%%%%%%%%%%%%%%%%%%%%%
\subsection{Minimal Polynomial}
\label{subsec:minimal-polynomial-computer-algebra}
%%%%%%%%%%%%%%%%%%%%%%%%%%%%%%%%%%%%%%%%%%%%%%%%%%%%%%%%%%%%

The minimal polynomial of algebraic number \(\alpha\) over \(\mathbb{Q}\)
is the monic irreducible polynomial of least degree satisfying

\[
\boxed{
m_\alpha(\alpha)=0.
}
\]

For

\[
\alpha=\sqrt{2},
\]

the minimal polynomial is

\[
\boxed{
m_\alpha(x)=x^2-2.
}
\]

%%%%%%%%%%%%%%%%%%%%%%%%%%%%%%%%%%%%%%%%%%%%%%%%%%%%%%%%%%%%
\subsection{Algebraic Number Fields}
\label{subsec:algebraic-number-fields-computation}
%%%%%%%%%%%%%%%%%%%%%%%%%%%%%%%%%%%%%%%%%%%%%%%%%%%%%%%%%%%%

A simple algebraic extension is written

\[
\boxed{
\mathbb{Q}(\alpha),
}
\]

where \(\alpha\) satisfies an irreducible polynomial of degree \(d\).

Every field element can be represented as

\[
\boxed{
a_0
+
a_1\alpha
+
\cdots
+
a_{d-1}\alpha^{d-1},
}
\]

with

\[
a_i\in\mathbb{Q}.
\]

Higher powers are reduced using the minimal polynomial.

%%%%%%%%%%%%%%%%%%%%%%%%%%%%%%%%%%%%%%%%%%%%%%%%%%%%%%%%%%%%
\subsection{Worked Example: Arithmetic in \(\mathbb{Q}(\sqrt{2})\)}
\label{subsec:worked-algebraic-number-arithmetic}
%%%%%%%%%%%%%%%%%%%%%%%%%%%%%%%%%%%%%%%%%%%%%%%%%%%%%%%%%%%%

\begin{workedexamplebox}{Reducing Powers}

Let

\[
\alpha=\sqrt{2}.
\]

Then

\[
\alpha^2=2.
\]

Compute

\[
(1+\alpha)^2.
\]

Expanding,

\[
1+2\alpha+\alpha^2.
\]

Since

\[
\alpha^2=2,
\]

we obtain

\begin{answerbox}
\[
\boxed{
(1+\sqrt{2})^2
=
3+2\sqrt{2}.
}
\]
\end{answerbox}

\end{workedexamplebox}

%%%%%%%%%%%%%%%%%%%%%%%%%%%%%%%%%%%%%%%%%%%%%%%%%%%%%%%%%%%%
\subsection{Polynomial Quotient Representation}
\label{subsec:polynomial-quotient-representation}
%%%%%%%%%%%%%%%%%%%%%%%%%%%%%%%%%%%%%%%%%%%%%%%%%%%%%%%%%%%%

If \(\alpha\) satisfies irreducible polynomial \(m(x)\), then

\[
\boxed{
\mathbb{Q}(\alpha)
\cong
\mathbb{Q}[x]/(m(x)).
}
\]

Arithmetic is performed modulo the relation

\[
\boxed{
m(x)=0.
}
\]

This quotient-ring viewpoint gives a clean computational representation.

%%%%%%%%%%%%%%%%%%%%%%%%%%%%%%%%%%%%%%%%%%%%%%%%%%%%%%%%%%%%
\subsection{Polynomial Ideals}
\label{subsec:computer-algebra-polynomial-ideals}
%%%%%%%%%%%%%%%%%%%%%%%%%%%%%%%%%%%%%%%%%%%%%%%%%%%%%%%%%%%%

Let

\[
R
=
F[x_1,\ldots,x_n].
\]

Given polynomials

\[
f_1,\ldots,f_m,
\]

their ideal is

\[
\boxed{
I
=
\langle
f_1,\ldots,f_m
\rangle.
}
\]

Every element of \(I\) has form

\[
\boxed{
h_1f_1+\cdots+h_mf_m,
}
\]

where

\[
h_i\in R.
\]

A central computational problem is determining whether a polynomial belongs
to \(I\).

%%%%%%%%%%%%%%%%%%%%%%%%%%%%%%%%%%%%%%%%%%%%%%%%%%%%%%%%%%%%
\subsection{Multivariate Polynomial Reduction}
\label{subsec:multivariate-polynomial-reduction}
%%%%%%%%%%%%%%%%%%%%%%%%%%%%%%%%%%%%%%%%%%%%%%%%%%%%%%%%%%%%

Given polynomials

\[
g_1,\ldots,g_s,
\]

one reduces a polynomial \(f\) by canceling leading monomials whenever

\[
\boxed{
\operatorname{LM}(g_i)
\mid
\operatorname{LM}(f).
}
\]

The process generalizes univariate polynomial division.

Unlike the univariate case, the remainder may depend on the reduction
order unless the divisors form a Gröbner basis.

%%%%%%%%%%%%%%%%%%%%%%%%%%%%%%%%%%%%%%%%%%%%%%%%%%%%%%%%%%%%
\subsection{Gröbner Basis}
\label{subsec:computer-algebra-groebner-basis}
%%%%%%%%%%%%%%%%%%%%%%%%%%%%%%%%%%%%%%%%%%%%%%%%%%%%%%%%%%%%

\begin{definitionbox}{Gröbner Basis}
Let

\[
I
\subseteq
F[x_1,\ldots,x_n]
\]

be an ideal with a fixed monomial ordering.

A finite set

\[
G
=
\{g_1,\ldots,g_s\}
\subseteq I
\]

is a Gröbner basis if the leading monomials of the \(g_i\) generate the
leading-term ideal of \(I\).
\end{definitionbox}

Equivalently, every nonzero polynomial in \(I\) has leading monomial
divisible by at least one

\[
\operatorname{LM}(g_i).
\]

%%%%%%%%%%%%%%%%%%%%%%%%%%%%%%%%%%%%%%%%%%%%%%%%%%%%%%%%%%%%
\subsection{Why Gröbner Bases Matter}
\label{subsec:why-groebner-bases}
%%%%%%%%%%%%%%%%%%%%%%%%%%%%%%%%%%%%%%%%%%%%%%%%%%%%%%%%%%%%

Gröbner bases provide algorithmic solutions to several problems:

\[
\boxed{
\text{ideal membership},
}
\]

\[
\boxed{
\text{polynomial-system solving},
}
\]

\[
\boxed{
\text{elimination},
}
\]

and

\[
\boxed{
\text{algebraic simplification modulo equations}.
}
\]

They play a role similar to row-reduced echelon form for systems of linear
equations.

%%%%%%%%%%%%%%%%%%%%%%%%%%%%%%%%%%%%%%%%%%%%%%%%%%%%%%%%%%%%
\subsection{Analogy With Gaussian Elimination}
\label{subsec:groebner-gaussian-analogy}
%%%%%%%%%%%%%%%%%%%%%%%%%%%%%%%%%%%%%%%%%%%%%%%%%%%%%%%%%%%%

Gaussian elimination transforms linear equations into a structured basis
that supports solving and consistency testing.

Gröbner-basis algorithms perform an analogous transformation for
polynomial ideals.

Thus one may think conceptually of

\[
\boxed{
\text{Gaussian elimination}
\leftrightarrow
\text{linear systems},
}
\]

and

\[
\boxed{
\text{Gröbner bases}
\leftrightarrow
\text{polynomial systems}.
}
\]

The polynomial problem is much more computationally difficult.

%%%%%%%%%%%%%%%%%%%%%%%%%%%%%%%%%%%%%%%%%%%%%%%%%%%%%%%%%%%%
\subsection{S-Polynomials}
\label{subsec:s-polynomials}
%%%%%%%%%%%%%%%%%%%%%%%%%%%%%%%%%%%%%%%%%%%%%%%%%%%%%%%%%%%%

Suppose

\[
f
\]

and

\[
g
\]

have leading terms

\[
\operatorname{LT}(f)
\]

and

\[
\operatorname{LT}(g).
\]

Let

\[
L
=
\operatorname{lcm}
(
\operatorname{LM}(f),
\operatorname{LM}(g)
).
\]

The S-polynomial is constructed to cancel the leading monomials:

\[
\boxed{
S(f,g)
=
\frac{
L
}{
\operatorname{LT}(f)
}
f
-
\frac{
L
}{
\operatorname{LT}(g)
}
g.
}
\]

S-polynomials reveal missing leading-term relations.

%%%%%%%%%%%%%%%%%%%%%%%%%%%%%%%%%%%%%%%%%%%%%%%%%%%%%%%%%%%%
\subsection{Buchberger's Algorithm}
\label{subsec:buchberger-algorithm}
%%%%%%%%%%%%%%%%%%%%%%%%%%%%%%%%%%%%%%%%%%%%%%%%%%%%%%%%%%%%

A basic Gröbner-basis algorithm is:

\begin{enumerate}
    \item begin with generators
    \[
    G=\{f_1,\ldots,f_m\};
    \]
    \item form S-polynomials of pairs in \(G\);
    \item reduce each S-polynomial by \(G\);
    \item if a nonzero remainder appears, add it to \(G\);
    \item repeat until every relevant S-polynomial reduces to zero.
\end{enumerate}

The final set is a Gröbner basis.

%%%%%%%%%%%%%%%%%%%%%%%%%%%%%%%%%%%%%%%%%%%%%%%%%%%%%%%%%%%%
\subsection{Buchberger Criterion}
\label{subsec:buchberger-criterion}
%%%%%%%%%%%%%%%%%%%%%%%%%%%%%%%%%%%%%%%%%%%%%%%%%%%%%%%%%%%%

A finite generating set \(G\) is a Gröbner basis if and only if every
S-polynomial of pairs of elements of \(G\) reduces to zero with respect to
\(G\).

This gives an algorithmic test for Gröbner-basis completeness.

%%%%%%%%%%%%%%%%%%%%%%%%%%%%%%%%%%%%%%%%%%%%%%%%%%%%%%%%%%%%
\subsection{Ideal Membership}
\label{subsec:ideal-membership}
%%%%%%%%%%%%%%%%%%%%%%%%%%%%%%%%%%%%%%%%%%%%%%%%%%%%%%%%%%%%

Suppose \(G\) is a Gröbner basis for ideal \(I\).

Reduce polynomial \(f\) by \(G\).

Then

\[
\boxed{
f\in I
}
\]

if and only if the Gröbner remainder is

\[
\boxed{
0.
}
\]

Thus ideal membership becomes a finite algorithmic calculation.

%%%%%%%%%%%%%%%%%%%%%%%%%%%%%%%%%%%%%%%%%%%%%%%%%%%%%%%%%%%%
\subsection{Worked Example: Ideal Membership Idea}
\label{subsec:worked-ideal-membership}
%%%%%%%%%%%%%%%%%%%%%%%%%%%%%%%%%%%%%%%%%%%%%%%%%%%%%%%%%%%%

\begin{workedexamplebox}{Simple Polynomial Ideal}

Consider the ideal

\[
I
=
\langle
x-y,
y-1
\rangle.
\]

The equations imply

\[
y=1,
\]

and

\[
x=1.
\]

Consider

\[
f=x+y-2.
\]

Substituting the ideal relations gives

\[
f
=
1+1-2
=
0.
\]

Therefore,

\begin{answerbox}
\[
\boxed{
x+y-2
\in
\langle
x-y,y-1
\rangle.
}
\]
\end{answerbox}

A Gröbner-basis reduction provides a systematic version of this reasoning.

\end{workedexamplebox}

%%%%%%%%%%%%%%%%%%%%%%%%%%%%%%%%%%%%%%%%%%%%%%%%%%%%%%%%%%%%
\subsection{Elimination Theorem}
\label{subsec:groebner-elimination-theorem}
%%%%%%%%%%%%%%%%%%%%%%%%%%%%%%%%%%%%%%%%%%%%%%%%%%%%%%%%%%%%

With lexicographic ordering, Gröbner bases can eliminate variables.

Suppose

\[
I
\subset
F[x_1,\ldots,x_n].
\]

A lexicographic Gröbner basis may contain polynomials involving only the
later variables.

This permits sequential solution of multivariate polynomial systems.

%%%%%%%%%%%%%%%%%%%%%%%%%%%%%%%%%%%%%%%%%%%%%%%%%%%%%%%%%%%%
\subsection{Triangular Polynomial Systems}
\label{subsec:triangular-polynomial-systems}
%%%%%%%%%%%%%%%%%%%%%%%%%%%%%%%%%%%%%%%%%%%%%%%%%%%%%%%%%%%%

An ideal may be transformed into equations such as

\[
\boxed{
g_1(x_n)=0,
}
\]

\[
\boxed{
g_2(x_{n-1},x_n)=0,
}
\]

and so on.

This resembles triangular linear systems.

Solve the univariate equation first, then substitute backward.

%%%%%%%%%%%%%%%%%%%%%%%%%%%%%%%%%%%%%%%%%%%%%%%%%%%%%%%%%%%%
\subsection{Zero-Dimensional Ideals Preview}
\label{subsec:zero-dimensional-ideals-preview}
%%%%%%%%%%%%%%%%%%%%%%%%%%%%%%%%%%%%%%%%%%%%%%%%%%%%%%%%%%%%

A polynomial system with finitely many solutions over an algebraic closure
is associated, under standard algebraic conditions, with a zero-dimensional
ideal.

For such systems, Gröbner bases can reduce solving to univariate
polynomials and back-substitution.

This provides an exact algebraic counterpart to numerical root finding.

%%%%%%%%%%%%%%%%%%%%%%%%%%%%%%%%%%%%%%%%%%%%%%%%%%%%%%%%%%%%
\subsection{Resultants}
\label{subsec:computer-algebra-resultants}
%%%%%%%%%%%%%%%%%%%%%%%%%%%%%%%%%%%%%%%%%%%%%%%%%%%%%%%%%%%%

For two univariate polynomials

\[
f(x)
\]

and

\[
g(x),
\]

their resultant satisfies

\[
\boxed{
\operatorname{Res}(f,g)=0
}
\]

if the polynomials share a common root over an algebraic closure, subject
to the usual nondegenerate interpretation.

Resultants can also eliminate a variable from multivariate polynomial
systems.

%%%%%%%%%%%%%%%%%%%%%%%%%%%%%%%%%%%%%%%%%%%%%%%%%%%%%%%%%%%%
\subsection{Sylvester Matrix}
\label{subsec:sylvester-matrix}
%%%%%%%%%%%%%%%%%%%%%%%%%%%%%%%%%%%%%%%%%%%%%%%%%%%%%%%%%%%%

The resultant of two univariate polynomials can be computed as the
determinant of their Sylvester matrix.

For

\[
f(x)
=
a_mx^m+\cdots+a_0
\]

and

\[
g(x)
=
b_nx^n+\cdots+b_0,
\]

the Sylvester matrix arranges shifted coefficient rows of \(f\) and \(g\).

Then

\[
\boxed{
\operatorname{Res}(f,g)
=
\det(S(f,g)).
}
\]

%%%%%%%%%%%%%%%%%%%%%%%%%%%%%%%%%%%%%%%%%%%%%%%%%%%%%%%%%%%%
\subsection{Worked Example: Simple Resultant}
\label{subsec:worked-resultant}
%%%%%%%%%%%%%%%%%%%%%%%%%%%%%%%%%%%%%%%%%%%%%%%%%%%%%%%%%%%%

\begin{workedexamplebox}{Two Linear Polynomials}

Let

\[
f(x)=ax+b,
\]

and

\[
g(x)=cx+d.
\]

The Sylvester matrix is

\[
\begin{pmatrix}
a&b\\
c&d
\end{pmatrix}.
\]

Therefore,

\[
\operatorname{Res}(f,g)
=
ad-bc.
\]

The polynomials have a common root precisely when

\[
ad-bc=0
\]

provided the equations are interpreted without degenerate zero-polynomial
cases.

\begin{answerbox}
\[
\boxed{
\operatorname{Res}(f,g)
=
ad-bc.
}
\]
\end{answerbox}

\end{workedexamplebox}

%%%%%%%%%%%%%%%%%%%%%%%%%%%%%%%%%%%%%%%%%%%%%%%%%%%%%%%%%%%%
\subsection{Exact Real Root Isolation}
\label{subsec:exact-real-root-isolation}
%%%%%%%%%%%%%%%%%%%%%%%%%%%%%%%%%%%%%%%%%%%%%%%%%%%%%%%%%%%%

A symbolic system may isolate real polynomial roots using rational
intervals.

For each real root \(\alpha\), find

\[
\boxed{
a<\alpha<b
}
\]

with rational endpoints such that the interval contains exactly one root.

This yields an exact representation even when no radical expression is
available.

%%%%%%%%%%%%%%%%%%%%%%%%%%%%%%%%%%%%%%%%%%%%%%%%%%%%%%%%%%%%
\subsection{Sturm Sequences Preview}
\label{subsec:sturm-sequences-preview}
%%%%%%%%%%%%%%%%%%%%%%%%%%%%%%%%%%%%%%%%%%%%%%%%%%%%%%%%%%%%

Sturm sequences provide an exact method for counting distinct real roots
of a polynomial inside an interval.

By repeatedly subdividing intervals, one can isolate all real roots.

This is an important example of exact real algebraic computation.

%%%%%%%%%%%%%%%%%%%%%%%%%%%%%%%%%%%%%%%%%%%%%%%%%%%%%%%%%%%%
\subsection{Symbolic Linear Algebra}
\label{subsec:computer-algebra-symbolic-linear-algebra}
%%%%%%%%%%%%%%%%%%%%%%%%%%%%%%%%%%%%%%%%%%%%%%%%%%%%%%%%%%%%

Computer algebra also performs matrix computations over exact domains.

Examples include:

\[
\boxed{
A\in\mathbb{Z}^{m\times n},
}
\]

\[
\boxed{
A\in\mathbb{Q}^{m\times n},
}
\]

or matrices over polynomial rings.

Typical tasks include:

\[
\boxed{
\text{rank},
}
\]

\[
\boxed{
\text{determinant},
}
\]

\[
\boxed{
\text{null space},
}
\]

and

\[
\boxed{
\text{linear-system solution}.
}
\]

%%%%%%%%%%%%%%%%%%%%%%%%%%%%%%%%%%%%%%%%%%%%%%%%%%%%%%%%%%%%
\subsection{Exact Gaussian Elimination}
\label{subsec:exact-gaussian-elimination}
%%%%%%%%%%%%%%%%%%%%%%%%%%%%%%%%%%%%%%%%%%%%%%%%%%%%%%%%%%%%

Gaussian elimination over \(\mathbb{Q}\) produces exact rational
solutions.

However, fractions may grow quickly.

For example, repeated pivot operations can create numerators and
denominators much larger than the original matrix entries.

This motivates fraction-free and modular approaches.

%%%%%%%%%%%%%%%%%%%%%%%%%%%%%%%%%%%%%%%%%%%%%%%%%%%%%%%%%%%%
\subsection{Fraction-Free Elimination}
\label{subsec:fraction-free-elimination}
%%%%%%%%%%%%%%%%%%%%%%%%%%%%%%%%%%%%%%%%%%%%%%%%%%%%%%%%%%%%

Fraction-free elimination reorganizes Gaussian elimination so divisions
are exact and fractions are postponed or avoided.

The intermediate values remain integers when the input matrix is integral.

This can dramatically reduce rational expression growth.

%%%%%%%%%%%%%%%%%%%%%%%%%%%%%%%%%%%%%%%%%%%%%%%%%%%%%%%%%%%%
\subsection{Bareiss-Type Elimination Preview}
\label{subsec:bareiss-elimination-preview}
%%%%%%%%%%%%%%%%%%%%%%%%%%%%%%%%%%%%%%%%%%%%%%%%%%%%%%%%%%%%

A fraction-free elimination scheme updates entries using divisions that
are known to be exact.

Conceptually, it preserves determinant-related scaling so intermediate
integers remain much smaller than those produced by naive rational
elimination.

This approach is useful for exact determinants and linear systems.

%%%%%%%%%%%%%%%%%%%%%%%%%%%%%%%%%%%%%%%%%%%%%%%%%%%%%%%%%%%%
\subsection{Exact Determinant Computation}
\label{subsec:exact-determinant}
%%%%%%%%%%%%%%%%%%%%%%%%%%%%%%%%%%%%%%%%%%%%%%%%%%%%%%%%%%%%

Possible strategies include:

\[
\boxed{
\text{fraction-free elimination},
}
\]

\[
\boxed{
\text{modular determinant computation},
}
\]

and

\[
\boxed{
\text{structured formulas}.
}
\]

Direct cofactor expansion is usually impractical for large matrices.

%%%%%%%%%%%%%%%%%%%%%%%%%%%%%%%%%%%%%%%%%%%%%%%%%%%%%%%%%%%%
\subsection{Modular Determinants}
\label{subsec:modular-determinants}
%%%%%%%%%%%%%%%%%%%%%%%%%%%%%%%%%%%%%%%%%%%%%%%%%%%%%%%%%%%%

For an integer matrix \(A\), compute

\[
\det(A)
\bmod p_i
\]

for several primes.

Use Chinese Remainder reconstruction to recover

\[
\boxed{
\det(A).
}
\]

If a bound on

\[
|\det(A)|
\]

is known, enough primes guarantee exact reconstruction.

%%%%%%%%%%%%%%%%%%%%%%%%%%%%%%%%%%%%%%%%%%%%%%%%%%%%%%%%%%%%
\subsection{Hadamard Bound}
\label{subsec:hadamard-bound-computer-algebra}
%%%%%%%%%%%%%%%%%%%%%%%%%%%%%%%%%%%%%%%%%%%%%%%%%%%%%%%%%%%%

For matrix \(A\) with row vectors \(r_i\),

\[
\boxed{
|\det(A)|
\leq
\prod_i
\|r_i\|_2.
}
\]

This provides an upper bound useful in determining how much modular
information is needed to reconstruct an exact determinant.

%%%%%%%%%%%%%%%%%%%%%%%%%%%%%%%%%%%%%%%%%%%%%%%%%%%%%%%%%%%%
\subsection{Hermite Normal Form}
\label{subsec:hermite-normal-form}
%%%%%%%%%%%%%%%%%%%%%%%%%%%%%%%%%%%%%%%%%%%%%%%%%%%%%%%%%%%%

The Hermite normal form gives a canonical triangular-like form for integer
matrices under multiplication by unimodular integer matrices.

It is useful for:

\[
\boxed{
\text{integer linear systems},
}
\]

\[
\boxed{
\text{lattices},
}
\]

and

\[
\boxed{
\text{module computations}.
}
\]

%%%%%%%%%%%%%%%%%%%%%%%%%%%%%%%%%%%%%%%%%%%%%%%%%%%%%%%%%%%%
\subsection{Smith Normal Form}
\label{subsec:smith-normal-form}
%%%%%%%%%%%%%%%%%%%%%%%%%%%%%%%%%%%%%%%%%%%%%%%%%%%%%%%%%%%%

For an integer matrix \(A\), the Smith normal form has diagonal structure

\[
\boxed{
UAV
=
D,
}
\]

where \(U\) and \(V\) are unimodular and

\[
\boxed{
D
=
\operatorname{diag}
(
d_1,\ldots,d_r,0,\ldots,0
),
}
\]

with divisibility conditions

\[
\boxed{
d_1\mid d_2\mid\cdots\mid d_r.
}
\]

This reveals structural information about finitely generated abelian
groups and integer modules.

%%%%%%%%%%%%%%%%%%%%%%%%%%%%%%%%%%%%%%%%%%%%%%%%%%%%%%%%%%%%
\subsection{Sparse Exact Linear Algebra}
\label{subsec:sparse-exact-linear-algebra}
%%%%%%%%%%%%%%%%%%%%%%%%%%%%%%%%%%%%%%%%%%%%%%%%%%%%%%%%%%%%

Exact matrices can also be sparse.

However, elimination may create fill-in.

Sparse exact algorithms therefore balance:

\[
\boxed{
\text{sparsity preservation}
}
\]

against

\[
\boxed{
\text{coefficient growth}.
}
\]

This is more difficult than ordinary sparse floating-point computation.

%%%%%%%%%%%%%%%%%%%%%%%%%%%%%%%%%%%%%%%%%%%%%%%%%%%%%%%%%%%%
\subsection{Expression Swell in Computer Algebra}
\label{subsec:computer-algebra-expression-swell}
%%%%%%%%%%%%%%%%%%%%%%%%%%%%%%%%%%%%%%%%%%%%%%%%%%%%%%%%%%%%

Exact algorithms may produce enormous intermediate expressions.

For example,

\[
\boxed{
\text{small input}
\longrightarrow
\text{huge intermediate result}
\longrightarrow
\text{small final answer}.
}
\]

This is expression swell.

It can dominate runtime and memory even when the theoretical number of
algebraic operations appears moderate.

%%%%%%%%%%%%%%%%%%%%%%%%%%%%%%%%%%%%%%%%%%%%%%%%%%%%%%%%%%%%
\subsection{Coefficient Growth}
\label{subsec:computer-algebra-coefficient-growth}
%%%%%%%%%%%%%%%%%%%%%%%%%%%%%%%%%%%%%%%%%%%%%%%%%%%%%%%%%%%%

Suppose initial coefficients contain at most \(b\) bits.

Intermediate coefficients may contain many more bits.

Thus the actual work depends on both:

\[
\boxed{
\text{number of algebraic operations}
}
\]

and

\[
\boxed{
\text{size of coefficients}.
}
\]

This motivates bit-complexity analysis.

%%%%%%%%%%%%%%%%%%%%%%%%%%%%%%%%%%%%%%%%%%%%%%%%%%%%%%%%%%%%
\subsection{Bit Complexity}
\label{subsec:computer-algebra-bit-complexity}
%%%%%%%%%%%%%%%%%%%%%%%%%%%%%%%%%%%%%%%%%%%%%%%%%%%%%%%%%%%%

In exact computation, multiplying two \(n\)-digit integers is more
expensive than multiplying machine-sized integers.

If \(M(n)\) denotes the cost of multiplying \(n\)-bit integers, then many
algebra algorithms are analyzed in terms of \(M(n)\).

Thus symbolic complexity often requires more detailed analysis than a
simple arithmetic-operation count.

%%%%%%%%%%%%%%%%%%%%%%%%%%%%%%%%%%%%%%%%%%%%%%%%%%%%%%%%%%%%
\subsection{Fast Integer Multiplication}
\label{subsec:fast-integer-multiplication}
%%%%%%%%%%%%%%%%%%%%%%%%%%%%%%%%%%%%%%%%%%%%%%%%%%%%%%%%%%%%

The elementary multiplication algorithm has quadratic digit complexity.

Faster algorithms use:

\[
\boxed{
\text{divide and conquer},
}
\]

and for sufficiently large integers,

\[
\boxed{
\text{transform-based multiplication}.
}
\]

These techniques are important because large-integer arithmetic appears
inside many exact algebra algorithms.

%%%%%%%%%%%%%%%%%%%%%%%%%%%%%%%%%%%%%%%%%%%%%%%%%%%%%%%%%%%%
\subsection{Evaluation and Interpolation Methods}
\label{subsec:evaluation-interpolation-computer-algebra}
%%%%%%%%%%%%%%%%%%%%%%%%%%%%%%%%%%%%%%%%%%%%%%%%%%%%%%%%%%%%

Some polynomial algorithms replace symbolic coefficient manipulation by:

\[
\boxed{
\text{evaluation at many points},
}
\]

followed by

\[
\boxed{
\text{pointwise computation},
}
\]

and then

\[
\boxed{
\text{interpolation}.
}
\]

This can accelerate multiplication, resultants, and other polynomial
operations.

%%%%%%%%%%%%%%%%%%%%%%%%%%%%%%%%%%%%%%%%%%%%%%%%%%%%%%%%%%%%
\subsection{Modular Images}
\label{subsec:modular-images}
%%%%%%%%%%%%%%%%%%%%%%%%%%%%%%%%%%%%%%%%%%%%%%%%%%%%%%%%%%%%

A large exact object may be mapped into smaller computational domains.

For example,

\[
A
\in
\mathbb{Z}^{n\times n}
\]

can be reduced to

\[
\boxed{
A\bmod p
\in
\mathbb{F}_p^{n\times n}.
}
\]

The modular image is easier to compute with.

Several images can later reconstruct the exact result.

%%%%%%%%%%%%%%%%%%%%%%%%%%%%%%%%%%%%%%%%%%%%%%%%%%%%%%%%%%%%
\subsection{Unlucky Primes}
\label{subsec:unlucky-primes}
%%%%%%%%%%%%%%%%%%%%%%%%%%%%%%%%%%%%%%%%%%%%%%%%%%%%%%%%%%%%

A prime used in modular computation can sometimes destroy important
structural information.

For example, a leading coefficient may become zero modulo \(p\), or rank
may decrease.

Such primes are often called unlucky primes.

Robust modular algorithms detect or avoid them.

%%%%%%%%%%%%%%%%%%%%%%%%%%%%%%%%%%%%%%%%%%%%%%%%%%%%%%%%%%%%
\subsection{Verification of Reconstructed Results}
\label{subsec:verification-reconstructed-results}
%%%%%%%%%%%%%%%%%%%%%%%%%%%%%%%%%%%%%%%%%%%%%%%%%%%%%%%%%%%%

After reconstructing an exact candidate, verify it against the original
problem.

For factorization,

\[
\boxed{
f
\stackrel{?}{=}
g_1g_2\cdots g_k.
}
\]

For a linear system,

\[
\boxed{
Ax
\stackrel{?}{=}
b.
}
\]

Exact verification is often inexpensive relative to discovering the
candidate result.

%%%%%%%%%%%%%%%%%%%%%%%%%%%%%%%%%%%%%%%%%%%%%%%%%%%%%%%%%%%%
\subsection{Randomized Computer Algebra}
\label{subsec:randomized-computer-algebra}
%%%%%%%%%%%%%%%%%%%%%%%%%%%%%%%%%%%%%%%%%%%%%%%%%%%%%%%%%%%%

Randomization can simplify exact algebraic algorithms.

Examples include choosing random:

\[
\boxed{
\text{evaluation points},
}
\]

\[
\boxed{
\text{primes},
}
\]

or

\[
\boxed{
\text{linear combinations}.
}
\]

Randomization may avoid pathological cases and reduce computational cost.

%%%%%%%%%%%%%%%%%%%%%%%%%%%%%%%%%%%%%%%%%%%%%%%%%%%%%%%%%%%%
\subsection{Las Vegas Algorithms}
\label{subsec:las-vegas-algebra}
%%%%%%%%%%%%%%%%%%%%%%%%%%%%%%%%%%%%%%%%%%%%%%%%%%%%%%%%%%%%

A Las Vegas algorithm always returns a correct result, but its runtime is
random.

For example, the algorithm may choose random parameters and restart if a
bad choice is detected.

Thus:

\[
\boxed{
\text{random runtime}
+
\text{guaranteed correctness}.
}
\]

%%%%%%%%%%%%%%%%%%%%%%%%%%%%%%%%%%%%%%%%%%%%%%%%%%%%%%%%%%%%
\subsection{Monte Carlo Algorithms in Algebra}
\label{subsec:monte-carlo-algebra}
%%%%%%%%%%%%%%%%%%%%%%%%%%%%%%%%%%%%%%%%%%%%%%%%%%%%%%%%%%%%

A Monte Carlo algorithm has bounded or controlled runtime but may return an
incorrect answer with small probability.

Repeated independent checks can reduce that probability.

Thus:

\[
\boxed{
\text{controlled runtime}
+
\text{small failure probability}.
}
\]

This terminology parallels randomized algorithms more broadly.

%%%%%%%%%%%%%%%%%%%%%%%%%%%%%%%%%%%%%%%%%%%%%%%%%%%%%%%%%%%%
\subsection{Polynomial Identity Testing}
\label{subsec:polynomial-identity-testing}
%%%%%%%%%%%%%%%%%%%%%%%%%%%%%%%%%%%%%%%%%%%%%%%%%%%%%%%%%%%%

Suppose one wants to test whether polynomial

\[
P(x_1,\ldots,x_n)
\]

is identically zero.

Expanding \(P\) may be expensive.

Instead, choose random input values and evaluate.

If

\[
P(a_1,\ldots,a_n)\neq0,
\]

then \(P\) is certainly nonzero.

If the evaluation is zero, repeat with new random values.

%%%%%%%%%%%%%%%%%%%%%%%%%%%%%%%%%%%%%%%%%%%%%%%%%%%%%%%%%%%%
\subsection{Schwartz--Zippel Principle Preview}
\label{subsec:schwartz-zippel-preview}
%%%%%%%%%%%%%%%%%%%%%%%%%%%%%%%%%%%%%%%%%%%%%%%%%%%%%%%%%%%%

For a nonzero polynomial of total degree \(d\), evaluating at random points
from a sufficiently large finite set yields zero with probability at most
approximately

\[
\boxed{
\frac{d}{|S|}.
}
\]

This provides a probabilistic identity test.

The error probability can be made very small by repeated trials.

%%%%%%%%%%%%%%%%%%%%%%%%%%%%%%%%%%%%%%%%%%%%%%%%%%%%%%%%%%%%
\subsection{Computer Algebra System Architecture}
\label{subsec:cas-architecture}
%%%%%%%%%%%%%%%%%%%%%%%%%%%%%%%%%%%%%%%%%%%%%%%%%%%%%%%%%%%%

A computer algebra system typically contains several layers.

Conceptually:

\[
\boxed{
\text{parser}
\rightarrow
\text{expression representation}
\rightarrow
\text{algebraic domains}
\rightarrow
\text{algorithms}
\rightarrow
\text{simplifier}
\rightarrow
\text{output}.
}
\]

Each layer has a distinct role.

%%%%%%%%%%%%%%%%%%%%%%%%%%%%%%%%%%%%%%%%%%%%%%%%%%%%%%%%%%%%
\subsection{Parser}
\label{subsec:cas-parser}
%%%%%%%%%%%%%%%%%%%%%%%%%%%%%%%%%%%%%%%%%%%%%%%%%%%%%%%%%%%%

The parser converts user notation into internal mathematical objects.

For example,

\[
\boxed{
x^2+2x+1
}
\]

must be recognized as a polynomial expression rather than an unstructured
text string.

Parsing determines operator precedence, function calls, and syntactic
structure.

%%%%%%%%%%%%%%%%%%%%%%%%%%%%%%%%%%%%%%%%%%%%%%%%%%%%%%%%%%%%
\subsection{Expression Kernel}
\label{subsec:expression-kernel}
%%%%%%%%%%%%%%%%%%%%%%%%%%%%%%%%%%%%%%%%%%%%%%%%%%%%%%%%%%%%

The expression kernel stores fundamental symbolic objects such as:

\[
\boxed{
\text{symbols},
}
\]

\[
\boxed{
\text{numbers},
}
\]

\[
\boxed{
\text{sums},
}
\]

\[
\boxed{
\text{products},
}
\]

and

\[
\boxed{
\text{powers}.
}
\]

Higher-level mathematics is built on these basic structures.

%%%%%%%%%%%%%%%%%%%%%%%%%%%%%%%%%%%%%%%%%%%%%%%%%%%%%%%%%%%%
\subsection{Domain System}
\label{subsec:cas-domain-system}
%%%%%%%%%%%%%%%%%%%%%%%%%%%%%%%%%%%%%%%%%%%%%%%%%%%%%%%%%%%%

A domain system tracks whether coefficients belong to:

\[
\boxed{
\mathbb{Z},
}
\]

\[
\boxed{
\mathbb{Q},
}
\]

\[
\boxed{
\mathbb{F}_p,
}
\]

or more complicated extensions.

Algorithms can then dispatch to methods suited to the domain.

For example, polynomial division differs over \(\mathbb{Z}\) and
\(\mathbb{Q}\).

%%%%%%%%%%%%%%%%%%%%%%%%%%%%%%%%%%%%%%%%%%%%%%%%%%%%%%%%%%%%
\subsection{Canonicalization}
\label{subsec:cas-canonicalization}
%%%%%%%%%%%%%%%%%%%%%%%%%%%%%%%%%%%%%%%%%%%%%%%%%%%%%%%%%%%%

Basic expressions are often automatically normalized.

Examples include:

\[
x+x
\longrightarrow
2x,
\]

\[
3+5
\longrightarrow
8,
\]

and reordering terms according to an internal canonical order.

Canonicalization makes expression comparison and hashing easier.

%%%%%%%%%%%%%%%%%%%%%%%%%%%%%%%%%%%%%%%%%%%%%%%%%%%%%%%%%%%%
\subsection{Simplification Engine}
\label{subsec:cas-simplification-engine}
%%%%%%%%%%%%%%%%%%%%%%%%%%%%%%%%%%%%%%%%%%%%%%%%%%%%%%%%%%%%

A simplifier applies transformations to obtain a preferred form.

Some simplifications are universally safe.

Others depend on:

\[
\boxed{
\text{domains},
}
\]

\[
\boxed{
\text{sign assumptions},
}
\]

or

\[
\boxed{
\text{branch conventions}.
}
\]

A reliable system must distinguish these cases.

%%%%%%%%%%%%%%%%%%%%%%%%%%%%%%%%%%%%%%%%%%%%%%%%%%%%%%%%%%%%
\subsection{Pattern Matcher}
\label{subsec:cas-pattern-matcher}
%%%%%%%%%%%%%%%%%%%%%%%%%%%%%%%%%%%%%%%%%%%%%%%%%%%%%%%%%%%%

Pattern matching enables rules such as

\[
\boxed{
a+a
\longrightarrow
2a,
}
\]

or

\[
\boxed{
\sin^2x+\cos^2x
\longrightarrow
1.
}
\]

Advanced matchers may support:

\[
\boxed{
\text{wildcards},
}
\]

\[
\boxed{
\text{sequence patterns},
}
\]

and

\[
\boxed{
\text{conditional rules}.
}
\]

%%%%%%%%%%%%%%%%%%%%%%%%%%%%%%%%%%%%%%%%%%%%%%%%%%%%%%%%%%%%
\subsection{Assumption Engine}
\label{subsec:cas-assumption-engine}
%%%%%%%%%%%%%%%%%%%%%%%%%%%%%%%%%%%%%%%%%%%%%%%%%%%%%%%%%%%%

Suppose the system knows

\[
\boxed{
x>0.
}
\]

Then it may simplify

\[
\sqrt{x^2}
\]

to

\[
\boxed{
x.
}
\]

Without that assumption, over the real numbers the safe simplification is

\[
\boxed{
|x|.
}
\]

Assumption systems therefore affect symbolic correctness.

%%%%%%%%%%%%%%%%%%%%%%%%%%%%%%%%%%%%%%%%%%%%%%%%%%%%%%%%%%%%
\subsection{Expression DAGs}
\label{subsec:cas-expression-dags}
%%%%%%%%%%%%%%%%%%%%%%%%%%%%%%%%%%%%%%%%%%%%%%%%%%%%%%%%%%%%

Repeated subexpressions can be shared.

Instead of representing

\[
f(g(x))+h(g(x))
\]

with two separate copies of \(g(x)\), a system may store one shared node.

This creates a directed acyclic graph.

Sharing reduces memory and supports faster equality checks.

%%%%%%%%%%%%%%%%%%%%%%%%%%%%%%%%%%%%%%%%%%%%%%%%%%%%%%%%%%%%
\subsection{Hashing Symbolic Objects}
\label{subsec:hashing-symbolic-objects}
%%%%%%%%%%%%%%%%%%%%%%%%%%%%%%%%%%%%%%%%%%%%%%%%%%%%%%%%%%%%

A normalized symbolic expression can be assigned a hash value.

Hashing supports:

\[
\boxed{
\text{fast lookup},
}
\]

\[
\boxed{
\text{memoization},
}
\]

and

\[
\boxed{
\text{common-subexpression detection}.
}
\]

Hash-based data structures are therefore important in large symbolic
computations.

%%%%%%%%%%%%%%%%%%%%%%%%%%%%%%%%%%%%%%%%%%%%%%%%%%%%%%%%%%%%
\subsection{Memoization}
\label{subsec:computer-algebra-memoization}
%%%%%%%%%%%%%%%%%%%%%%%%%%%%%%%%%%%%%%%%%%%%%%%%%%%%%%%%%%%%

If a symbolic operation is expensive and repeated, store the result.

For function \(F\),

\[
\boxed{
x
\mapsto
F(x)
}
\]

can be cached.

Future calls with the same normalized input reuse the stored answer.

This is especially effective for recursive symbolic algorithms.

%%%%%%%%%%%%%%%%%%%%%%%%%%%%%%%%%%%%%%%%%%%%%%%%%%%%%%%%%%%%
\subsection{Lazy Algebraic Objects}
\label{subsec:lazy-algebraic-objects}
%%%%%%%%%%%%%%%%%%%%%%%%%%%%%%%%%%%%%%%%%%%%%%%%%%%%%%%%%%%%

Some exact objects can be represented without fully expanding them.

For example,

\[
\boxed{
(x+1)^{1000}
}
\]

can remain as a power.

Only if coefficients are explicitly requested must the polynomial be
expanded.

Lazy representation helps control expression swell.

%%%%%%%%%%%%%%%%%%%%%%%%%%%%%%%%%%%%%%%%%%%%%%%%%%%%%%%%%%%%
\subsection{Black-Box Polynomials}
\label{subsec:black-box-polynomials}
%%%%%%%%%%%%%%%%%%%%%%%%%%%%%%%%%%%%%%%%%%%%%%%%%%%%%%%%%%%%

A polynomial may be accessible only through evaluation:

\[
\boxed{
x
\mapsto
f(x).
}
\]

The coefficient representation may be unknown or too large to store.

Black-box algorithms operate using evaluations rather than explicit
coefficients.

This is useful in interpolation, identity testing, and sparse polynomial
reconstruction.

%%%%%%%%%%%%%%%%%%%%%%%%%%%%%%%%%%%%%%%%%%%%%%%%%%%%%%%%%%%%
\subsection{Sparse Interpolation Preview}
\label{subsec:sparse-interpolation-preview}
%%%%%%%%%%%%%%%%%%%%%%%%%%%%%%%%%%%%%%%%%%%%%%%%%%%%%%%%%%%%

Suppose a polynomial has enormous degree but only a few nonzero terms.

Dense interpolation would be wasteful.

Sparse interpolation attempts to recover only the nonzero monomials and
their coefficients from carefully chosen evaluations.

This is an important example of adapting algorithms to algebraic
structure.

%%%%%%%%%%%%%%%%%%%%%%%%%%%%%%%%%%%%%%%%%%%%%%%%%%%%%%%%%%%%
\subsection{Symbolic-Numeric Algebra}
\label{subsec:symbolic-numeric-algebra}
%%%%%%%%%%%%%%%%%%%%%%%%%%%%%%%%%%%%%%%%%%%%%%%%%%%%%%%%%%%%

Exact algebraic algorithms may become prohibitively expensive.

Numerical methods can provide approximate candidates.

Then symbolic computation can reconstruct or verify exact structure.

A common strategy is:

\[
\boxed{
\text{numerical approximation}
\rightarrow
\text{exact reconstruction}
\rightarrow
\text{verification}.
}
\]

%%%%%%%%%%%%%%%%%%%%%%%%%%%%%%%%%%%%%%%%%%%%%%%%%%%%%%%%%%%%
\subsection{Integer Relation Detection Preview}
\label{subsec:integer-relation-detection}
%%%%%%%%%%%%%%%%%%%%%%%%%%%%%%%%%%%%%%%%%%%%%%%%%%%%%%%%%%%%

Suppose a high-precision numerical constant appears to satisfy a relation

\[
\boxed{
a_0
+
a_1\alpha_1
+
\cdots
+
a_n\alpha_n
=
0,
}
\]

with small integers \(a_i\).

Integer-relation algorithms search for such coefficient vectors.

This can help identify exact formulas from high-precision numerical data.

Any conjectured relation must then be verified independently.

%%%%%%%%%%%%%%%%%%%%%%%%%%%%%%%%%%%%%%%%%%%%%%%%%%%%%%%%%%%%
\subsection{Exact Reconstruction From Floating Point}
\label{subsec:exact-reconstruction-floating}
%%%%%%%%%%%%%%%%%%%%%%%%%%%%%%%%%%%%%%%%%%%%%%%%%%%%%%%%%%%%

Suppose a numerical computation suggests

\[
x
\approx
0.333333333333.
\]

One may conjecture

\[
\boxed{
x=\frac13.
}
\]

Rational approximation algorithms can search for small numerator and
denominator pairs.

However, approximate agreement alone is not proof.

Exact substitution or mathematical verification is required.

%%%%%%%%%%%%%%%%%%%%%%%%%%%%%%%%%%%%%%%%%%%%%%%%%%%%%%%%%%%%
\subsection{Algebraic Complexity}
\label{subsec:algebraic-complexity}
%%%%%%%%%%%%%%%%%%%%%%%%%%%%%%%%%%%%%%%%%%%%%%%%%%%%%%%%%%%%

Computer algebra complexity depends on several dimensions:

\[
\boxed{
\text{degree},
}
\]

\[
\boxed{
\text{number of variables},
}
\]

\[
\boxed{
\text{coefficient bit size},
}
\]

and

\[
\boxed{
\text{number of terms}.
}
\]

A method efficient in one parameter may still become infeasible because
another parameter grows.

%%%%%%%%%%%%%%%%%%%%%%%%%%%%%%%%%%%%%%%%%%%%%%%%%%%%%%%%%%%%
\subsection{Gröbner-Basis Complexity}
\label{subsec:groebner-complexity}
%%%%%%%%%%%%%%%%%%%%%%%%%%%%%%%%%%%%%%%%%%%%%%%%%%%%%%%%%%%%

Gröbner-basis computation can be extremely expensive in the worst case.

Intermediate polynomial degrees and term counts can grow dramatically.

Practical performance depends strongly on:

\[
\boxed{
\text{monomial ordering},
}
\]

\[
\boxed{
\text{system structure},
}
\]

and

\[
\boxed{
\text{implementation strategy}.
}
\]

Thus a mathematically equivalent reformulation can change runtime by orders
of magnitude.

%%%%%%%%%%%%%%%%%%%%%%%%%%%%%%%%%%%%%%%%%%%%%%%%%%%%%%%%%%%%
\subsection{Choosing Monomial Orders}
\label{subsec:choosing-monomial-orders}
%%%%%%%%%%%%%%%%%%%%%%%%%%%%%%%%%%%%%%%%%%%%%%%%%%%%%%%%%%%%

Lexicographic order is useful for elimination and triangular solving.

Graded reverse lexicographic order is often computationally cheaper for
basis construction.

A common strategy is:

\[
\boxed{
\text{compute in an efficient order}
}
\]

then

\[
\boxed{
\text{convert to an elimination-friendly order}.
}
\]

Order choice is part of algorithm design.

%%%%%%%%%%%%%%%%%%%%%%%%%%%%%%%%%%%%%%%%%%%%%%%%%%%%%%%%%%%%
\subsection{Computer Algebra Workflow}
\label{subsec:computer-algebra-workflow}
%%%%%%%%%%%%%%%%%%%%%%%%%%%%%%%%%%%%%%%%%%%%%%%%%%%%%%%%%%%%

A practical computer-algebra workflow is:

\begin{enumerate}
    \item identify the exact algebraic domain;
    \item preserve integers and rationals exactly;
    \item choose dense or sparse representations;
    \item choose monomial ordering for multivariate polynomials;
    \item normalize input objects;
    \item remove trivial common factors;
    \item use primitive parts when working over the integers;
    \item detect repeated factors;
    \item choose direct, modular, or evaluation-based algorithms;
    \item control intermediate coefficient growth;
    \item exploit sparsity;
    \item use finite-field images when useful;
    \item reconstruct exact results;
    \item verify reconstructed candidates;
    \item avoid unnecessary expression expansion;
    \item use Gröbner bases only when polynomial-ideal structure is needed;
    \item choose elimination order carefully;
    \item use symbolic-numeric hybrids when exact computation is too large;
    \item document coefficient domains and assumptions;
    \item independently verify important results.
\end{enumerate}

%%%%%%%%%%%%%%%%%%%%%%%%%%%%%%%%%%%%%%%%%%%%%%%%%%%%%%%%%%%%
\subsection{Choosing a Computer Algebra Strategy}
\label{subsec:choosing-computer-algebra-strategy}
%%%%%%%%%%%%%%%%%%%%%%%%%%%%%%%%%%%%%%%%%%%%%%%%%%%%%%%%%%%%

For small univariate polynomials, direct exact algorithms are often
appropriate.

For large integer coefficients, consider:

\[
\boxed{
\text{modular methods}.
}
\]

For multivariate polynomial ideals, consider:

\[
\boxed{
\text{Gröbner bases}.
}
\]

For variable elimination, consider:

\[
\boxed{
\text{resultants}
}
\]

or

\[
\boxed{
\text{lexicographic Gröbner bases}.
}
\]

For huge sparse polynomials, consider:

\[
\boxed{
\text{sparse and black-box methods}.
}
\]

For large exact matrices, consider:

\[
\boxed{
\text{fraction-free or modular linear algebra}.
}
\]

%%%%%%%%%%%%%%%%%%%%%%%%%%%%%%%%%%%%%%%%%%%%%%%%%%%%%%%%%%%%
\subsection{Common Errors}
\label{subsec:computer-algebra-common-errors}
%%%%%%%%%%%%%%%%%%%%%%%%%%%%%%%%%%%%%%%%%%%%%%%%%%%%%%%%%%%%

\subsubsection{Ignoring the Coefficient Domain}

The polynomial

\[
x^2-2
\]

is irreducible over

\[
\mathbb{Q}
\]

but reducible over

\[
\mathbb{R}.
\]

Factorization questions are meaningless without specifying the domain.

%%%%%%%%%%%%%%%%%%%%%%%%%%%%%%%%%%%%%%%%%%%%%%%%%%%%%%%%%%%%
\subsubsection{Using Dense Representation for Extremely Sparse Polynomials}

A polynomial like

\[
x^{1000000}+x+1
\]

should not require storage for one million zero coefficients.

Choose a representation suited to sparsity.

%%%%%%%%%%%%%%%%%%%%%%%%%%%%%%%%%%%%%%%%%%%%%%%%%%%%%%%%%%%%
\subsubsection{Using Sparse Representation for Dense Low-Degree Polynomials}

Sparse structures add indexing overhead.

Dense storage is often simpler and faster when most coefficients are
nonzero.

%%%%%%%%%%%%%%%%%%%%%%%%%%%%%%%%%%%%%%%%%%%%%%%%%%%%%%%%%%%%
\subsubsection{Performing Rational Arithmetic Without Reduction}

Repeated fraction operations can produce unnecessarily huge numerators and
denominators.

Normalize fractions systematically.

%%%%%%%%%%%%%%%%%%%%%%%%%%%%%%%%%%%%%%%%%%%%%%%%%%%%%%%%%%%%
\subsubsection{Expanding Before Factorization}

Full expansion can cause expression swell.

Preserve useful structure unless expansion is required by the algorithm.

%%%%%%%%%%%%%%%%%%%%%%%%%%%%%%%%%%%%%%%%%%%%%%%%%%%%%%%%%%%%
\subsubsection{Using Ordinary Division Over a Nonfield}

Leading coefficients may not be invertible in rings such as

\[
\mathbb{Z}.
\]

Pseudo-division or domain conversion may be required.

%%%%%%%%%%%%%%%%%%%%%%%%%%%%%%%%%%%%%%%%%%%%%%%%%%%%%%%%%%%%
\subsubsection{Ignoring Primitive Parts}

For integer polynomials, separating content from primitive part often
reduces unnecessary coefficient growth.

%%%%%%%%%%%%%%%%%%%%%%%%%%%%%%%%%%%%%%%%%%%%%%%%%%%%%%%%%%%%
\subsubsection{Ignoring Repeated Factors Before Factorization}

Square-free factorization is an important preprocessing step.

Repeated factors can complicate later algorithms.

%%%%%%%%%%%%%%%%%%%%%%%%%%%%%%%%%%%%%%%%%%%%%%%%%%%%%%%%%%%%
\subsubsection{Using a Bad Modular Prime}

An unlucky prime can change degree, rank, or factor structure.

Modular algorithms should detect inconsistent images.

%%%%%%%%%%%%%%%%%%%%%%%%%%%%%%%%%%%%%%%%%%%%%%%%%%%%%%%%%%%%
\subsubsection{Reconstructing Before Enough Modular Information Is Available}

Chinese Remainder reconstruction gives a result modulo \(M\).

The modulus must be large enough relative to the possible exact answer.

Bounds or verification are required.

%%%%%%%%%%%%%%%%%%%%%%%%%%%%%%%%%%%%%%%%%%%%%%%%%%%%%%%%%%%%
\subsubsection{Assuming a Numerical Guess Is Exact}

A floating-point number that looks rational or algebraic may merely be
close.

Exact reconstruction must be verified.

%%%%%%%%%%%%%%%%%%%%%%%%%%%%%%%%%%%%%%%%%%%%%%%%%%%%%%%%%%%%
\subsubsection{Using Gröbner Bases Without Choosing an Appropriate Monomial Order}

Different monomial orders can produce dramatically different computation
times and basis structures.

The desired output should guide the order.

%%%%%%%%%%%%%%%%%%%%%%%%%%%%%%%%%%%%%%%%%%%%%%%%%%%%%%%%%%%%
\subsubsection{Assuming Multivariate Division Has a Unique Remainder}

For arbitrary divisors, multivariate remainder can depend on reduction
order.

Uniqueness is obtained in the appropriate Gröbner-basis setting with fixed
monomial order and reduction convention.

%%%%%%%%%%%%%%%%%%%%%%%%%%%%%%%%%%%%%%%%%%%%%%%%%%%%%%%%%%%%
\subsubsection{Assuming Every Gröbner Basis Is Reduced}

A Gröbner basis need not be minimal or reduced.

A reduced Gröbner basis has additional normalization properties and is
unique for a fixed monomial order over a field.

%%%%%%%%%%%%%%%%%%%%%%%%%%%%%%%%%%%%%%%%%%%%%%%%%%%%%%%%%%%%
\subsubsection{Using Lexicographic Order for Every Gröbner Computation}

Lex order is valuable for elimination but can be much more expensive than
graded orders.

A two-stage order strategy may be better.

%%%%%%%%%%%%%%%%%%%%%%%%%%%%%%%%%%%%%%%%%%%%%%%%%%%%%%%%%%%%
\subsubsection{Using Cofactor Expansion for Large Determinants}

Cofactor expansion grows combinatorially.

Use elimination, fraction-free, or modular determinant algorithms.

%%%%%%%%%%%%%%%%%%%%%%%%%%%%%%%%%%%%%%%%%%%%%%%%%%%%%%%%%%%%
\subsubsection{Using Exact Rational Gaussian Elimination Naively}

Intermediate fractions can become enormous.

Fraction-free or modular methods often perform much better.

%%%%%%%%%%%%%%%%%%%%%%%%%%%%%%%%%%%%%%%%%%%%%%%%%%%%%%%%%%%%
\subsubsection{Ignoring Bit Complexity}

An algorithm performing few operations on huge integers may be slower than
one performing more operations on small modular values.

Coefficient size matters.

%%%%%%%%%%%%%%%%%%%%%%%%%%%%%%%%%%%%%%%%%%%%%%%%%%%%%%%%%%%%
\subsubsection{Failing to Verify Randomized Algebra Results}

Randomized methods may rely on probabilistic assumptions or random choices.

Whenever possible, perform an exact final verification.

%%%%%%%%%%%%%%%%%%%%%%%%%%%%%%%%%%%%%%%%%%%%%%%%%%%%%%%%%%%%
\subsubsection{Confusing Computer Algebra With Numerical Algebra}

An exact determinant computation and a floating-point determinant
computation solve related but different computational problems.

Their algorithms, errors, and complexity considerations differ.

%%%%%%%%%%%%%%%%%%%%%%%%%%%%%%%%%%%%%%%%%%%%%%%%%%%%%%%%%%%%
\subsection{Applications}
\label{subsec:computer-algebra-applications}
%%%%%%%%%%%%%%%%%%%%%%%%%%%%%%%%%%%%%%%%%%%%%%%%%%%%%%%%%%%%

\begin{applicationbox}

\textbf{Abstract Algebra.}
Computer algebra performs exact ring, field, polynomial, ideal, and module
computations.

\medskip

\textbf{Number Theory.}
Large-integer arithmetic, modular computation, factorization, algebraic
numbers, and exact matrix normal forms are central computational tools.

\medskip

\textbf{Algebraic Geometry.}
Gröbner bases, ideals, elimination, resultants, and polynomial systems are
fundamental computational methods.

\medskip

\textbf{Cryptography.}
Finite fields, polynomial arithmetic, modular arithmetic, and algebraic
structures support many cryptographic algorithms.

\medskip

\textbf{Coding Theory.}
Polynomial and finite-field calculations are used to construct and decode
error-correcting codes.

\medskip

\textbf{Optimization.}
Polynomial optimization and exact KKT elimination can use algebraic
methods.

\medskip

\textbf{Geometry.}
Exact intersections of algebraic curves and surfaces can be reduced to
polynomial-system calculations.

\medskip

\textbf{Scientific Computing.}
Computer algebra derives exact formulas, Jacobians, reduced models, and
code that can later be evaluated numerically.

\end{applicationbox}

%%%%%%%%%%%%%%%%%%%%%%%%%%%%%%%%%%%%%%%%%%%%%%%%%%%%%%%%%%%%
\subsection{Exercises}
\label{subsec:computer-algebra-exercises}
%%%%%%%%%%%%%%%%%%%%%%%%%%%%%%%%%%%%%%%%%%%%%%%%%%%%%%%%%%%%

\begin{exercise}[1]
Define computer algebra.
\end{exercise}

\begin{exercise}[1]
Distinguish symbolic computation from computer algebra.
\end{exercise}

\begin{exercise}[1]
Define a polynomial ring.
\end{exercise}

\begin{exercise}[1]
Define a monomial order.
\end{exercise}

\begin{exercise}[1]
Define the leading monomial of a polynomial.
\end{exercise}

\begin{exercise}[1]
Define a polynomial ideal.
\end{exercise}

\begin{exercise}[1]
Define a Gröbner basis.
\end{exercise}

\begin{exercise}[1]
Define an S-polynomial.
\end{exercise}

\begin{exercise}[1]
Define a resultant conceptually.
\end{exercise}

\begin{exercise}[1]
Define an algebraic number.
\end{exercise}

\begin{exercise}[2]
Represent

\[
3x^5-2x+7
\]

using both dense and sparse representations.
\end{exercise}

\begin{exercise}[2]
Determine which representation is better for

\[
x^{10000}+x^2+1.
\]
\end{exercise}

\begin{exercise}[2]
Add two polynomials using coefficient arrays.
\end{exercise}

\begin{exercise}[2]
Multiply two small polynomials by coefficient convolution.
\end{exercise}

\begin{exercise}[3]
Compare the naive polynomial multiplication cost with FFT-based
multiplication conceptually.
\end{exercise}

\begin{exercise}[2]
Perform polynomial division over

\[
\mathbb{Q}[x].
\]
\end{exercise}

\begin{exercise}[3]
Explain why ordinary polynomial division may fail over

\[
\mathbb{Z}[x].
\]
\end{exercise}

\begin{exercise}[3]
Perform a simple pseudo-division.
\end{exercise}

\begin{exercise}[2]
Compute a polynomial gcd using the Euclidean algorithm.
\end{exercise}

\begin{exercise}[3]
Find the content and primitive part of an integer polynomial.
\end{exercise}

\begin{exercise}[3]
Use

\[
\gcd(f,f')
\]

to detect repeated factors.
\end{exercise}

\begin{exercise}[3]
Construct a square-free decomposition of a simple polynomial.
\end{exercise}

\begin{exercise}[3]
Factor a polynomial modulo a small prime.
\end{exercise}

\begin{exercise}[3]
Explain the role of finite-field factorization in integer polynomial
factorization.
\end{exercise}

\begin{exercise}[3]
Explain Hensel lifting conceptually.
\end{exercise}

\begin{exercise}[2]
Solve a system of congruences using the Chinese Remainder Theorem.
\end{exercise}

\begin{exercise}[3]
Explain how modular images can reduce coefficient growth.
\end{exercise}

\begin{exercise}[3]
Explain rational reconstruction.
\end{exercise}

\begin{exercise}[2]
Find the minimal polynomial of

\[
\sqrt{3}.
\]
\end{exercise}

\begin{exercise}[3]
Reduce

\[
(1+\sqrt{3})^3
\]

to the basis

\[
1,\sqrt{3}.
\]
\end{exercise}

\begin{exercise}[3]
Explain why

\[
\mathbb{Q}(\alpha)
\cong
\mathbb{Q}[x]/(m_\alpha(x)).
\]
\end{exercise}

\begin{exercise}[2]
Compare monomials under lexicographic ordering.
\end{exercise}

\begin{exercise}[2]
Compare monomials under graded lexicographic ordering.
\end{exercise}

\begin{exercise}[2]
Compare monomials under graded reverse lexicographic ordering.
\end{exercise}

\begin{exercise}[3]
Find the leading monomial and leading term of a multivariate polynomial
under a specified monomial order.
\end{exercise}

\begin{exercise}[3]
Perform one multivariate polynomial reduction step.
\end{exercise}

\begin{exercise}[3]
Explain why multivariate division can depend on divisor order.
\end{exercise}

\begin{exercise}[3]
Construct an S-polynomial for two given polynomials.
\end{exercise}

\begin{exercise}[3]
Reduce an S-polynomial by a given basis.
\end{exercise}

\begin{exercise}[3]
Explain Buchberger's algorithm.
\end{exercise}

\begin{exercise}[3]
Test ideal membership using a supplied Gröbner basis.
\end{exercise}

\begin{exercise}[3]
Explain why lexicographic Gröbner bases support elimination.
\end{exercise}

\begin{exercise}[3]
Solve a small triangular polynomial system.
\end{exercise}

\begin{exercise}[3]
Explain the analogy between Gaussian elimination and Gröbner bases.
\end{exercise}

\begin{exercise}[2]
Construct the Sylvester matrix of two linear polynomials.
\end{exercise}

\begin{exercise}[3]
Compute their resultant.
\end{exercise}

\begin{exercise}[3]
Explain why a zero resultant indicates a common root under suitable
conditions.
\end{exercise}

\begin{exercise}[3]
Describe exact real root isolation.
\end{exercise}

\begin{exercise}[3]
Explain the purpose of a Sturm sequence.
\end{exercise}

\begin{exercise}[2]
Perform exact Gaussian elimination on a rational matrix.
\end{exercise}

\begin{exercise}[3]
Show how intermediate rational coefficients grow.
\end{exercise}

\begin{exercise}[3]
Explain fraction-free elimination.
\end{exercise}

\begin{exercise}[3]
Compute a determinant modulo several primes.
\end{exercise}

\begin{exercise}[3]
Reconstruct an integer determinant using the Chinese Remainder Theorem.
\end{exercise}

\begin{exercise}[3]
Use the Hadamard bound to estimate a determinant magnitude.
\end{exercise}

\begin{exercise}[3]
Explain Hermite normal form.
\end{exercise}

\begin{exercise}[3]
Explain Smith normal form.
\end{exercise}

\begin{exercise}[3]
Describe the divisibility condition on Smith normal form diagonal entries.
\end{exercise}

\begin{exercise}[3]
Explain why exact sparse elimination may create fill-in.
\end{exercise}

\begin{exercise}[3]
Give an example of coefficient growth in an exact computation.
\end{exercise}

\begin{exercise}[3]
Explain bit complexity.
\end{exercise}

\begin{exercise}[3]
Explain why large-integer multiplication matters in computer algebra.
\end{exercise}

\begin{exercise}[3]
Describe the evaluation-interpolation strategy for polynomial
multiplication.
\end{exercise}

\begin{exercise}[3]
Explain what an unlucky prime is.
\end{exercise}

\begin{exercise}[3]
Design a verification step after modular reconstruction.
\end{exercise}

\begin{exercise}[3]
Distinguish Las Vegas and Monte Carlo randomized algorithms.
\end{exercise}

\begin{exercise}[3]
Describe randomized polynomial identity testing.
\end{exercise}

\begin{exercise}[3]
Explain the Schwartz--Zippel principle conceptually.
\end{exercise}

\begin{exercise}[3]
Draw or describe the architecture of a computer algebra system.
\end{exercise}

\begin{exercise}[3]
Explain the role of the expression kernel.
\end{exercise}

\begin{exercise}[3]
Explain the role of the algebraic-domain system.
\end{exercise}

\begin{exercise}[3]
Give examples of automatic canonicalization.
\end{exercise}

\begin{exercise}[3]
Explain why symbolic assumptions affect simplification.
\end{exercise}

\begin{exercise}[3]
Explain the benefit of expression DAGs.
\end{exercise}

\begin{exercise}[3]
Explain memoization in a symbolic algorithm.
\end{exercise}

\begin{exercise}[3]
Explain lazy representation of

\[
(x+1)^{1000}.
\]
\end{exercise}

\begin{exercise}[3]
Describe a black-box polynomial.
\end{exercise}

\begin{exercise}[3]
Explain sparse polynomial interpolation conceptually.
\end{exercise}

\begin{exercise}[3]
Explain symbolic-numeric reconstruction.
\end{exercise}

\begin{exercise}[3]
Explain why a conjectured rational reconstruction requires verification.
\end{exercise}

\begin{exercise}[4]
Study Karatsuba integer multiplication.
\end{exercise}

\begin{exercise}[4]
Study transform-based integer multiplication.
\end{exercise}

\begin{exercise}[4]
Study fast polynomial multiplication.
\end{exercise}

\begin{exercise}[4]
Study subresultant polynomial remainder sequences.
\end{exercise}

\begin{exercise}[4]
Study modular gcd algorithms.
\end{exercise}

\begin{exercise}[4]
Study finite-field polynomial factorization algorithms.
\end{exercise}

\begin{exercise}[4]
Study Hensel lifting in detail.
\end{exercise}

\begin{exercise}[4]
Study rational reconstruction algorithms.
\end{exercise}

\begin{exercise}[4]
Study exact algebraic-number arithmetic.
\end{exercise}

\begin{exercise}[4]
Study primitive element representations of number fields.
\end{exercise}

\begin{exercise}[4]
Derive Buchberger's criterion.
\end{exercise}

\begin{exercise}[4]
Implement Buchberger's algorithm conceptually for a small polynomial
system.
\end{exercise}

\begin{exercise}[4]
Study Gröbner-basis complexity.
\end{exercise}

\begin{exercise}[4]
Study elimination ideals.
\end{exercise}

\begin{exercise}[4]
Study change-of-order algorithms for Gröbner bases.
\end{exercise}

\begin{exercise}[4]
Study resultant computation by subresultants.
\end{exercise}

\begin{exercise}[4]
Study Sturm root-counting algorithms.
\end{exercise}

\begin{exercise}[4]
Study fraction-free Gaussian elimination.
\end{exercise}

\begin{exercise}[4]
Study modular exact linear algebra.
\end{exercise}

\begin{exercise}[4]
Study Hermite normal form algorithms.
\end{exercise}

\begin{exercise}[4]
Study Smith normal form algorithms.
\end{exercise}

\begin{exercise}[4]
Study sparse exact matrix algorithms.
\end{exercise}

\begin{exercise}[4]
Study randomized rank and determinant algorithms.
\end{exercise}

\begin{exercise}[4]
Study polynomial identity testing.
\end{exercise}

\begin{exercise}[4]
Study sparse interpolation algorithms.
\end{exercise}

\begin{exercise}[4]
Study integer-relation algorithms.
\end{exercise}

\begin{exercise}[4]
Investigate the design of a small computer algebra kernel.
\end{exercise}

%%%%%%%%%%%%%%%%%%%%%%%%%%%%%%%%%%%%%%%%%%%%%%%%%%%%%%%%%%%%
\subsection{Conceptual Summary}
\label{subsec:computer-algebra-summary}
%%%%%%%%%%%%%%%%%%%%%%%%%%%%%%%%%%%%%%%%%%%%%%%%%%%%%%%%%%%%

Computer algebra develops exact algorithms for structured mathematical
objects.

Important coefficient domains include

\[
\boxed{
\mathbb{Z},
\quad
\mathbb{Q},
\quad
\mathbb{F}_p,
}
\]

and polynomial extensions.

Polynomial computation includes:

\[
\boxed{
\text{addition}
+
\text{multiplication}
+
\text{division}
+
\text{gcd}
+
\text{factorization}.
}
\]

Modular methods use

\[
\boxed{
\text{finite-field images}
\rightarrow
\text{exact reconstruction}.
}
\]

Gröbner bases provide computational tools for

\[
\boxed{
\text{ideals}
+
\text{elimination}
+
\text{multivariate polynomial systems}.
}
\]

Exact linear algebra uses

\[
\boxed{
\text{fraction-free elimination}
+
\text{modular methods}
+
\text{normal forms}.
}
\]

Efficient computer algebra must control

\[
\boxed{
\text{expression swell}
+
\text{coefficient growth}
+
\text{bit complexity}.
}
\]

Modern systems combine deterministic, modular, randomized, symbolic, and
numerical techniques.

%%%%%%%%%%%%%%%%%%%%%%%%%%%%%%%%%%%%%%%%%%%%%%%%%%%%%%%%%%%%
\subsection{Section Connections}
\label{subsec:computer-algebra-connections}
%%%%%%%%%%%%%%%%%%%%%%%%%%%%%%%%%%%%%%%%%%%%%%%%%%%%%%%%%%%%

\chapterconnection{
Computer Algebra provides the algorithmic foundation behind the symbolic
operations introduced in Section~\ref{sec:symbolic-computation}.
Polynomial arithmetic, modular computation, algebraic-number
representation, Gröbner bases, resultants, exact linear algebra, and
symbolic system architecture connect directly with the algebra and number
theory developed earlier in the text. These algorithms also support exact
preprocessing for numerical optimization, differential equations,
algebraic geometry, coding theory, cryptography, and scientific computing.
The next section, Mathematical Simulation, combines mathematical models
with numerical algorithms to construct computational experiments for
dynamic, stochastic, physical, biological, and engineering systems.
}

%%%%%%%%%%%%%%%%%%%%%%%%%%%%%%%%%%%%%%%%%%%%%%%%%%%%%%%%%%%%
% SECTION SOURCE TRACKING
%%%%%%%%%%%%%%%%%%%%%%%%%%%%%%%%%%%%%%%%%%%%%%%%%%%%%%%%%%%%

%------------------------------------------------------------
% 11.9 Computer Algebra
%------------------------------------------------------------
%
% Source basis:
%   - Standard computer algebra algorithms.
%   - Standard exact polynomial and integer computation.
%   - Standard Gröbner-basis, elimination, modular, and exact linear
%     algebra methods.
%
% Used for:
%   - computer algebra definition;
%   - symbolic versus computer algebra distinction;
%   - exact algebraic domains;
%   - arbitrary-size integers and rationals;
%   - normalization;
%   - polynomial rings;
%   - dense and sparse polynomial representations;
%   - multivariate monomials;
%   - monomial orderings;
%   - lexicographic, graded lexicographic, and graded reverse
%     lexicographic orders;
%   - leading monomial and leading term;
%   - polynomial addition and multiplication;
%   - fast polynomial multiplication;
%   - polynomial division and pseudo-division;
%   - polynomial gcd;
%   - primitive parts and content;
%   - Gauss's lemma computationally;
%   - subresultant PRS;
%   - square-free factorization;
%   - finite-field factorization;
%   - modular polynomial factorization;
%   - Hensel lifting preview;
%   - Chinese Remainder reconstruction;
%   - rational reconstruction;
%   - modular algorithms;
%   - algebraic numbers;
%   - minimal polynomials;
%   - algebraic number fields;
%   - quotient representations;
%   - polynomial ideals;
%   - multivariate reduction;
%   - Gröbner bases;
%   - S-polynomials;
%   - Buchberger's algorithm and criterion;
%   - ideal membership;
%   - elimination;
%   - zero-dimensional systems preview;
%   - resultants and Sylvester matrices;
%   - exact real-root isolation;
%   - Sturm sequence preview;
%   - symbolic linear algebra;
%   - fraction-free elimination;
%   - exact and modular determinants;
%   - Hadamard bound;
%   - Hermite normal form;
%   - Smith normal form;
%   - sparse exact linear algebra;
%   - expression swell and coefficient growth;
%   - bit complexity;
%   - fast integer arithmetic;
%   - evaluation and interpolation;
%   - modular images and unlucky primes;
%   - exact verification;
%   - randomized computer algebra;
%   - Las Vegas and Monte Carlo algebra algorithms;
%   - polynomial identity testing;
%   - Schwartz--Zippel preview;
%   - computer algebra system architecture;
%   - parsers, expression kernels, domains, simplifiers, assumptions,
%     pattern matching, hashing, DAGs, memoization, and lazy evaluation;
%   - black-box polynomials;
%   - sparse interpolation preview;
%   - symbolic-numeric algebra;
%   - integer-relation and reconstruction previews;
%   - applications and exercises.
%
% Notes:
%   - Mathematical Simulation is developed next in Section 11.10.
%   - High-Performance Mathematical Computing follows in Section 11.11.
%   - Symbolic Computation appears in Section 11.8.
%
%%%%%%%%%%%%%%%%%%%%%%%%%%%%%%%%%%%%%%%%%%%%%%%%%%%%%%%%%%%%